%% Volume I, Chapter 13: Vectors and Matrices
\chapter{Vectors and Matrices}
\label{ch:vectors-matrices}

\abstract{Almost every problem of scientific computing ends as a computation with matrices: a system of linear equations, a least squares fit, an eigenvalue problem. A matrix looks like a table of numbers, but it is better understood as the record of a linear map, and the product of two matrices is the composition of their maps. From this single fact the rules of matrix algebra follow, including the surprising ones: the order of the factors matters, and two nonzero matrices can have the product zero. The same product can be read entry by entry, by columns, by rows, as a sum of outer products or block by block, and each reading is the starting point of algorithms in later volumes.}

\Cref{sec:vectors-rn} introduced vectors in $\mathbb R^n$, the inner product, matrices and their products, but only as far as the calculus of several variables needed them, and several facts were stated there with a reference to the present part of the book. This chapter begins Part~III, on linear algebra, and builds the language of vectors and matrices systematically. We keep every convention of \cref{ch:several-variables}: vectors are bold lower-case letters and are columns when they meet matrices, matrices are italic capitals with entries $a_{ij}$, the inner product is $\langle\mathbf x,\mathbf y\rangle=\mathbf x^{\mathsf T}\mathbf y$, and $\|\mathbf x\|$ is the Euclidean norm. Facts that \cref{sec:vectors-rn} proved are quoted with a reference, and those it only stated are proved here, unless they belong to a later chapter.

The guiding idea is that a matrix is the record of a linear map. The columns of the matrix are the images of the standard basis vectors, the product of a matrix with a vector is a combination of the columns, and the product of two matrices is the matrix of the composed map. Almost every rule of the chapter is a consequence of this idea, and so are the exceptions to the familiar rules of arithmetic: composition of maps is not commutative, so neither is the product of matrices; a map that squeezes the plane onto a line cannot be undone, so not every nonzero matrix has an inverse. The second theme is that the same product can be organised in several ways, by entries, by columns, by rows, as a sum of outer products, or in blocks, and that these readings, which give the same numbers in exact arithmetic, lead to different algorithms. The dense linear algebra of Volume~VIII, the sparse matrices of Volume~IX and the array products of Volume~VI are all built on them.

The chapter is about the vocabulary and the arithmetic of vectors and matrices, the words that the rest of Part~III uses. The questions that need more theory are postponed: whether a linear system has a solution and how to find it is the subject of \cref{ch:linear-systems}, the abstract notions of vector space, basis and dimension come in \cref{ch:vector-spaces}, determinants in \cref{ch:determinants}, and orthogonal projections onto subspaces and the QR factorisation in \cref{ch:inner-products}. Entries are real numbers unless stated otherwise; \cref{sec:complex-matrices} explains what changes for complex entries, which \cref{ch:eigenvalues} and later chapters need.

The chapter is organised as follows. \Cref{sec:vectors-in-rn} reviews the arithmetic of vectors and introduces linear combinations, spans, lines and planes. \Cref{sec:inner-product-length} treats the inner product, lengths, angles, orthogonality and the norms $\|\cdot\|_1$ and $\|\cdot\|_\infty$. \Cref{sec:matrices-linear-maps} establishes the correspondence between matrices and linear maps and collects examples, and \cref{sec:matrix-multiplication} derives the matrix product from the composition of maps and studies its four readings, its powers and its cost. \Cref{sec:transpose,sec:inverse} treat the transpose and the inverse, \cref{sec:special-matrices} the special kinds of matrices that occur again and again: diagonal, triangular, permutation, orthogonal, banded and others. \Cref{sec:block-matrices} computes with blocks and derives the Schur complement and the Sherman--Morrison--Woodbury formula, \cref{sec:trace-norms} introduces the trace and the norms of matrices and proves that the inverse of a matrix survives small perturbations, and \cref{sec:complex-matrices} extends everything to complex entries. \Cref{sec:computing-matrices} counts the operations of the basic computations and measures how the time of a matrix product grows with its size. As in the preceding chapters this is a review; but nearly all the results of this chapter have short proofs, and we prove them all, except the theorem that a one-sided inverse of a square matrix is two-sided, which \cref{ch:linear-systems} proves.

\section{Vectors in \texorpdfstring{$\mathbb R^n$}{R\^{}n}}
\label{sec:vectors-in-rn}

\subsection{The Arithmetic of Vectors}

A vector $\mathbf x=(x_1,\ldots,x_n)\in\mathbb R^n$ can be thought of in three ways, and all three are useful. It is a list of $n$ numbers: the measurements of $n$ quantities, the prices of $n$ goods, the coefficients of a polynomial, the grey values of the $n$ pixels of an image. It is a point of $n$-dimensional space, with coordinates $x_1,\ldots,x_n$. And it is an arrow, a displacement from the origin to that point, or from any point $\mathbf p$ to $\mathbf p+\mathbf x$. The arithmetic of \cref{sec:vectors-rn}, addition and multiplication by scalars component by component, serves all three views at once. In the first view, adding the price lists of two shops or converting all prices into another currency are vector operations. In the third, $\mathbf x+\mathbf y$ is the displacement $\mathbf x$ followed by the displacement $\mathbf y$, and $t\mathbf x$ is the displacement $\mathbf x$ stretched by the factor $t$.

\begin{proposition}[Rules of vector arithmetic]
\label{prop:vector-rules}
For all $\mathbf x,\mathbf y,\mathbf z\in\mathbb R^n$ and $s,t\in\mathbb R$:
\begin{enumerate}
\item $(\mathbf x+\mathbf y)+\mathbf z=\mathbf x+(\mathbf y+\mathbf z)$;
\item $\mathbf x+\mathbf y=\mathbf y+\mathbf x$;
\item $\mathbf x+\mathbf 0=\mathbf x$;
\item $\mathbf x+(-\mathbf x)=\mathbf 0$;
\item $s(t\mathbf x)=(st)\mathbf x$;
\item $1\mathbf x=\mathbf x$;
\item $t(\mathbf x+\mathbf y)=t\mathbf x+t\mathbf y$;
\item $(s+t)\mathbf x=s\mathbf x+t\mathbf x$.
\end{enumerate}
\end{proposition}

\begin{proof}
Each rule is an identity between two vectors, and two vectors are equal when their components are equal. The $i$th components of the two sides of (i) are $(x_i+y_i)+z_i$ and $x_i+(y_i+z_i)$, which agree by the associative law of the real numbers. In the same way, (ii) to (viii) reduce in each component to the commutative law of addition, $x_i+0=x_i$, $x_i+(-x_i)=0$, the associative law of multiplication, $1\cdot x_i=x_i$, and the two distributive laws $t(x_i+y_i)=tx_i+ty_i$ and $(s+t)x_i=sx_i+tx_i$ (\cref{sec:reals}).
\end{proof}

The eight rules are worth listing because they are all that most arguments about vectors use. Whenever a set carries an addition and a multiplication by scalars that satisfy them, the arguments go through unchanged. This is the case for the polynomials, for the continuous functions on an interval and for the matrices of a fixed shape, and it is the reason why \cref{ch:vector-spaces} takes the eight rules as the axioms of an abstract \term{vector space}. Other familiar rules follow from the eight. For instance, $0\mathbf x=\mathbf 0$ for every $\mathbf x$: by (viii), $0\mathbf x=(0+0)\mathbf x=0\mathbf x+0\mathbf x$, and adding $-(0\mathbf x)$ to both sides and using (i), (iii) and (iv) leaves $\mathbf 0=0\mathbf x$. In $\mathbb R^n$ this is evident from the components, but the derivation from the rules shows that it holds in every vector space. Similarly $t\mathbf 0=\mathbf 0$, $(-1)\mathbf x=-\mathbf x$, and $t\mathbf x=\mathbf 0$ only if $t=0$ or $\mathbf x=\mathbf 0$.

\subsection{Linear Combinations and Span}

The two operations combine into the most important construction of linear algebra.

\begin{definition}[Linear combination and span]
\label{def:span}
Let $\mathbf v_1,\ldots,\mathbf v_k\in\mathbb R^n$. A \term{linear combination} of these vectors is a vector
\[
t_1\mathbf v_1+t_2\mathbf v_2+\cdots+t_k\mathbf v_k=\sum_{j=1}^kt_j\mathbf v_j\qquad(t_1,\ldots,t_k\in\mathbb R),
\]
and the numbers $t_j$ are its \term{coefficients}. The set of all linear combinations is the \term{span} of the vectors,
\[
\operatorname{span}\{\mathbf v_1,\ldots,\mathbf v_k\}\coloneqq\Bigl\{\sum_{j=1}^kt_j\mathbf v_j: t_1,\ldots,t_k\in\mathbb R\Bigr\},
\]
and the vectors are said to \term{span} this set. The span of the empty family is $\{\mathbf 0\}$.
\end{definition}

The span of a single vector $\mathbf v\ne\mathbf 0$ is the line $\{t\mathbf v: t\in\mathbb R\}$ through the origin in the direction of $\mathbf v$. In the plane, the span of $(1,2)$ is the line $y=2x$, and the span of $(1,2)$ and $(2,4)$ is the same line, because $(2,4)=2(1,2)$ adds nothing new: every combination $s(1,2)+t(2,4)=(s+2t)(1,2)$ is already a multiple of $(1,2)$. The span of $(2,1)$ and $(1,2)$, by contrast, is the whole plane. \Cref{fig:linear-combinations} shows why: the combinations $s\mathbf u+t\mathbf v$ with integer coefficients form a slanted grid, and every point of the plane lies in one of its cells, so that it can be reached by moving along $\mathbf u$ and then along $\mathbf v$. For instance, $(3,0)=2(2,1)-(1,2)$. In general, the coefficients of a vector $(x_1,x_2)$ are found by solving the equations $2s+t=x_1$, $s+2t=x_2$, a linear system, whose solution $s=(2x_1-x_2)/3$, $t=(2x_2-x_1)/3$ exists for every $(x_1,x_2)$; \cref{ch:linear-systems} studies such systems in general. The standard basis vectors span $\mathbb R^n$, since every vector is the combination $\mathbf x=\sum_ix_i\mathbf e_i$ of them, with its components as coefficients.

\begin{figure}
\sdpanel{A}{\begin{tikzpicture}[font=\scriptsize, scale=0.78]
\fill[sdfill] (0,0) -- (2,1) -- (3,3) -- (1,2) -- cycle;
\draw[line width=0.4pt, -{Latex[length=1.6mm]}] (-0.4,0) -- (4.1,0) node[below] {$x_1$};
\draw[line width=0.4pt, -{Latex[length=1.6mm]}] (0,-0.4) -- (0,3.7) node[left] {$x_2$};
\draw[black!60, dashed, line width=0.5pt] (2,1) -- (3,3) -- (1,2);
\draw[line width=0.8pt, -{Latex[length=2mm]}] (0,0) -- (2,1);
\draw[line width=0.8pt, -{Latex[length=2mm]}] (0,0) -- (1,2);
\draw[line width=1.3pt, -{Latex[length=2mm]}] (0,0) -- (3,3);
\node[below right] at (2,1) {$\mathbf u=(2,1)$};
\node[above left] at (1,2) {$\mathbf v=(1,2)$};
\node[above] at (3,3) {$\mathbf u+\mathbf v=(3,3)$};
\end{tikzpicture}}\hfill
\sdpanel{B}{\begin{tikzpicture}[font=\scriptsize, scale=0.6]
\begin{scope}
\clip (-2.2,-2.2) rectangle (5.4,4.3);
\foreach \s in {-6,...,6} {\draw[black!30, line width=0.35pt] ({2*\s-6},{\s-12}) -- ({2*\s+6},{\s+12});}
\foreach \t in {-6,...,6} {\draw[black!30, line width=0.35pt] ({\t-12},{2*\t-6}) -- ({\t+12},{2*\t+6});}
\end{scope}
\draw[line width=0.4pt, -{Latex[length=1.6mm]}] (-2.2,0) -- (5.6,0) node[below] {$x_1$};
\draw[line width=0.4pt, -{Latex[length=1.6mm]}] (0,-2.2) -- (0,4.5) node[left] {$x_2$};
\draw[line width=0.8pt, -{Latex[length=2mm]}] (0,0) -- (2,1);
\draw[line width=0.8pt, -{Latex[length=2mm]}] (0,0) -- (1,2);
\draw[line width=1.1pt, -{Latex[length=2mm]}] (2,1) -- (4,2);
\draw[line width=1.1pt, -{Latex[length=2mm]}] (4,2) -- (3,0);
\fill (3,0) circle (2.2pt);
\node[below right] at (1.75,0.95) {$\mathbf u$};
\node[left] at (0.5,1.2) {$\mathbf v$};
\node[above] at (4,2.05) {$2\mathbf u$};
\node[below] at (3,-0.05) {$2\mathbf u-\mathbf v$};
\end{tikzpicture}}
\caption{Linear combinations of $\mathbf u=(2,1)$ and $\mathbf v=(1,2)$. (A) The sum $\mathbf u+\mathbf v$ is the diagonal of the parallelogram spanned by $\mathbf u$ and $\mathbf v$ (shaded). (B) The lines $\{s\mathbf u+t\mathbf v: t\in\mathbb R\}$ for integer $s$ and $\{s\mathbf u+t\mathbf v: s\in\mathbb R\}$ for integer $t$ form a slanted grid whose intersections are the combinations with integer coefficients. Every point of the plane lies in a cell of the grid and is a combination $s\mathbf u+t\mathbf v$; the point $(3,0)$, for instance, is reached by going twice along $\mathbf u$ and once against $\mathbf v$.}
\label{fig:linear-combinations}
\end{figure}

In space, the span of two vectors that are not multiples of each other is a plane through the origin. For $\mathbf v_1=(1,0,1)$ and $\mathbf v_2=(0,1,1)$, the combinations are $s\mathbf v_1+t\mathbf v_2=(s,t,s+t)$, so every point of the span satisfies $z=x+y$; conversely, every point $(x,y,x+y)$ of this plane is the combination $x\mathbf v_1+y\mathbf v_2$. The span is therefore exactly the plane $z=x+y$. A third vector enlarges the span to all of $\mathbb R^3$ if and only if it does not lie in this plane, a statement that \cref{ch:vector-spaces} makes precise with the notions of linear independence and dimension.

\begin{proposition}[Spans]
\label{prop:span}
Let $S=\operatorname{span}\{\mathbf v_1,\ldots,\mathbf v_k\}$. Then $S$ contains $\mathbf 0$ and every $\mathbf v_j$, and it contains $\mathbf x+\mathbf y$ and $t\mathbf x$ whenever it contains $\mathbf x$ and $\mathbf y$. Every set with these two properties that contains $\mathbf v_1,\ldots,\mathbf v_k$ contains $S$.
\end{proposition}

\begin{proof}
$\mathbf 0$ is the combination with all coefficients $0$, and $\mathbf v_j$ the one with $t_j=1$ and all other coefficients $0$. If $\mathbf x=\sum_js_j\mathbf v_j$ and $\mathbf y=\sum_jt_j\mathbf v_j$, then $\mathbf x+\mathbf y=\sum_j(s_j+t_j)\mathbf v_j$ and $t\mathbf x=\sum_j(ts_j)\mathbf v_j$, by the rules of \cref{prop:vector-rules}. A set $T$ that contains the $\mathbf v_j$ and is closed under addition and multiplication by scalars contains every $t_j\mathbf v_j$ and then, by induction on the number of terms, every sum $\sum_jt_j\mathbf v_j$.
\end{proof}

So the span is the smallest set containing the given vectors that is closed under the two operations of vector arithmetic. Such closed sets are the \term{subspaces} of \cref{ch:vector-spaces}.

\subsection{Lines and Planes}

Lines and planes that do not pass through the origin are spans shifted away from it. The \term{line} through a point $\mathbf p$ in the direction $\mathbf v\ne\mathbf 0$ is the set
\[
\{\mathbf p+t\mathbf v: t\in\mathbb R\},
\]
and the map $t\mapsto\mathbf p+t\mathbf v$ is a \term{parametrisation} of it, a path in the sense of \cref{def:path} if $t$ is restricted to an interval. The line through two different points $\mathbf p$ and $\mathbf q$ has the direction $\mathbf q-\mathbf p$, so its points are
\[
\mathbf p+t(\mathbf q-\mathbf p)=(1-t)\mathbf p+t\mathbf q\qquad(t\in\mathbb R).
\]
The parameter $t=0$ gives $\mathbf p$, $t=1$ gives $\mathbf q$, $t=1/2$ the midpoint $(\mathbf p+\mathbf q)/2$, and the values $0\le t\le1$ the \term{segment} from $\mathbf p$ to $\mathbf q$. The line through $(1,2,0)$ and $(3,1,4)$, for instance, consists of the points $(1+2t,2-t,4t)$; it meets the plane $z=2$ at $t=1/2$, in the point $(2,3/2,2)$. The combinations $(1-t)\mathbf p+t\mathbf q$ with $0\le t\le1$, whose coefficients are nonnegative and add up to $1$, are the \term{convex combinations} of $\mathbf p$ and $\mathbf q$, and \cref{def:convex-rn} called a set convex if it contains the segment between any two of its points.

Similarly, the \term{plane} through $\mathbf p$ spanned by two vectors $\mathbf u$ and $\mathbf v$, neither of which is a multiple of the other, is the set
\[
\{\mathbf p+s\mathbf u+t\mathbf v: s,t\in\mathbb R\}.
\]
The plane through the three points $\mathbf e_1$, $\mathbf e_2$ and $\mathbf e_3$ of $\mathbb R^3$ is obtained with $\mathbf p=\mathbf e_1$, $\mathbf u=\mathbf e_2-\mathbf e_1$ and $\mathbf v=\mathbf e_3-\mathbf e_1$; its points $(1-s-t,s,t)$ have coordinates that add up to $1$, and every point with this property is of the form $(1-s-t,s,t)$. The plane is therefore the set $x+y+z=1$. Its points with nonnegative coordinates form the triangle with the vertices $\mathbf e_1,\mathbf e_2,\mathbf e_3$, the set of convex combinations of the three points.

The two descriptions of the plane, by a parametrisation and by an equation, are typical. A parametrisation produces the points of the set; an equation tests whether a given point belongs to it. Passing from one description to the other means solving linear equations, as in the example, and \cref{sec:inner-product-length} gives the equation of a hyperplane a geometric meaning. In $\mathbb R^n$ the same construction with $k$ vectors gives a \term{flat} of dimension $k$, a shifted span, provided that the $k$ vectors are linearly independent, a notion that \cref{ch:vector-spaces} defines.

\section{The Inner Product, Length, and Angle}
\label{sec:inner-product-length}

\subsection{The Inner Product and the Euclidean Norm}

The inner product $\langle\mathbf x,\mathbf y\rangle=\sum_ix_iy_i$ and the Euclidean norm $\|\mathbf x\|=\langle\mathbf x,\mathbf x\rangle^{1/2}$ were defined in \cref{def:inner-product}. \Cref{prop:inner-product} showed that the inner product is symmetric, bilinear and positive, and \cref{prop:norm} that the norm is positive, homogeneous and satisfies the triangle inequality, with the Cauchy--Schwarz inequality $|\langle\mathbf x,\mathbf y\rangle|\le\|\mathbf x\|\|\mathbf y\|$ (\cref{thm:cauchy-schwarz-vectors}) as the main step. In the language of matrices the inner product is the product of a row and a column, $\langle\mathbf x,\mathbf y\rangle=\mathbf x^{\mathsf T}\mathbf y$, and in that of lists it is the total that one computes from quantities and prices: if $\mathbf x$ lists the numbers of items bought and $\mathbf y$ their prices, $\langle\mathbf x,\mathbf y\rangle$ is the bill.

The Cauchy--Schwarz inequality is so central that a second proof is instructive. \Cref{thm:cauchy-schwarz} derived it by minimising a sum of squares over a parameter; the following identity, due to Lagrange, writes the gap between the two sides as a sum of squares directly~\cite[Chap.~1]{steele2004}.

\begin{proposition}[Lagrange's identity]
\label{prop:lagrange-identity}
For all $\mathbf x,\mathbf y\in\mathbb R^n$,
\[
\|\mathbf x\|^2\|\mathbf y\|^2-\langle\mathbf x,\mathbf y\rangle^2=\sum_{1\le i<j\le n}\bigl(x_iy_j-x_jy_i\bigr)^2 .
\]
Consequently $|\langle\mathbf x,\mathbf y\rangle|\le\|\mathbf x\|\|\mathbf y\|$, with equality if and only if one of the two vectors is a multiple of the other.
\end{proposition}

\begin{proof}
Multiplying out the sums,
\[
\|\mathbf x\|^2\|\mathbf y\|^2-\langle\mathbf x,\mathbf y\rangle^2=\sum_{i=1}^n\sum_{j=1}^nx_i^2y_j^2-\sum_{i=1}^n\sum_{j=1}^nx_iy_ix_jy_j .
\]
Each double sum is unchanged when the names $i$ and $j$ are exchanged, so the difference is half the sum of the two versions:
\[
\frac12\sum_{i=1}^n\sum_{j=1}^n\bigl(x_i^2y_j^2+x_j^2y_i^2-2x_iy_jx_jy_i\bigr)=\frac12\sum_{i=1}^n\sum_{j=1}^n\bigl(x_iy_j-x_jy_i\bigr)^2 .
\]
The terms with $i=j$ vanish, and the terms for $(i,j)$ and $(j,i)$ are equal, which gives the identity. The right side is nonnegative, which is the inequality, and it vanishes exactly when $x_iy_j=x_jy_i$ for all $i$ and $j$. If $\mathbf x=\mathbf 0$, then $\mathbf x=0\mathbf y$. Otherwise some $x_i\ne0$, and with $t\coloneqq y_i/x_i$ the conditions give $y_j=tx_j$ for all $j$, that is, $\mathbf y=t\mathbf x$. Conversely, if $\mathbf y=t\mathbf x$ or $\mathbf x=t\mathbf y$, every term $x_iy_j-x_jy_i$ is zero.
\end{proof}

For $n=3$ the three terms of the sum are the squares of the components of the cross product $\mathbf x\times\mathbf y$, and the identity becomes $\|\mathbf x\times\mathbf y\|^2=\|\mathbf x\|^2\|\mathbf y\|^2-\langle\mathbf x,\mathbf y\rangle^2$, part~(iii) of \cref{prop:cross-product}. For $n=2$ there is a single term, $(x_1y_2-x_2y_1)^2$, the square of the determinant of the matrix with the columns $\mathbf x$ and $\mathbf y$, that is, of the area of the parallelogram they span. In every dimension, the sum on the right is the square of the area of that parallelogram, and $\|\mathbf x\|^2\|\mathbf y\|^2-\langle\mathbf x,\mathbf y\rangle^2$ is the \term{Gram determinant} of the two vectors, which \cref{ch:determinants} generalises to $k$ vectors.

\subsection{Angles and Orthogonality}

By the Cauchy--Schwarz inequality, the \term{angle} $\theta\in[0,\pi]$ between nonzero vectors, defined by $\cos\theta=\langle\mathbf x,\mathbf y\rangle/(\|\mathbf x\|\|\mathbf y\|)$ in \eqref{eq:angle}, is well defined, and the law of cosines justifies the name. For $\mathbf x=(1,2,2)$ and $\mathbf y=(2,-1,2)$, both of length $3$, the inner product is $2-2+4=4$, so $\cos\theta=4/9$ and $\theta=1.1102\ldots$, about $63.6$ degrees. The sign of the inner product says whether the angle is acute ($\langle\mathbf x,\mathbf y\rangle>0$), right ($=0$) or obtuse ($<0$).

Vectors with $\langle\mathbf x,\mathbf y\rangle=0$ are \term{orthogonal}, $\mathbf x\perp\mathbf y$. The zero vector is orthogonal to every vector, and it is the only vector orthogonal to itself, by \cref{prop:inner-product}\,(iii). A vector of length $1$ is a \term{unit vector}, and every nonzero vector $\mathbf x$ is a positive multiple of exactly one unit vector, its direction $\mathbf x/\|\mathbf x\|$; passing from $\mathbf x$ to $\mathbf x/\|\mathbf x\|$ is called \term{normalising} $\mathbf x$. Vectors $\mathbf q_1,\ldots,\mathbf q_k$ are \term{orthonormal} if they are unit vectors and pairwise orthogonal, $\langle\mathbf q_i,\mathbf q_j\rangle=\delta_{ij}$, where the \term{Kronecker delta} $\delta_{ij}$ is $1$ for $i=j$ and $0$ otherwise. The standard basis vectors are orthonormal, and so are $(1,1)/\sqrt2$ and $(1,-1)/\sqrt2$ in the plane.

\begin{proposition}[Pythagoras for several vectors]
\label{prop:pythagoras}
If $\mathbf v_1,\ldots,\mathbf v_k$ are pairwise orthogonal, then $\|\mathbf v_1+\cdots+\mathbf v_k\|^2=\|\mathbf v_1\|^2+\cdots+\|\mathbf v_k\|^2$. In particular, if $\mathbf q_1,\ldots,\mathbf q_k$ are orthonormal, then $\bigl\|\sum_jt_j\mathbf q_j\bigr\|^2=\sum_jt_j^2$.
\end{proposition}

\begin{proof}
By bilinearity, $\bigl\|\sum_i\mathbf v_i\bigr\|^2=\bigl\langle\sum_i\mathbf v_i,\sum_j\mathbf v_j\bigr\rangle=\sum_{i,j}\langle\mathbf v_i,\mathbf v_j\rangle$, and only the terms with $i=j$ are nonzero. For the second statement apply the first to $\mathbf v_j=t_j\mathbf q_j$, whose squared length is $t_j^2$.
\end{proof}

For an orthonormal family the coefficients of a combination are easy to recover: taking the inner product of $\mathbf x=\sum_jt_j\mathbf q_j$ with $\mathbf q_i$ gives $\langle\mathbf q_i,\mathbf x\rangle=\sum_jt_j\delta_{ij}=t_i$. With respect to the orthonormal vectors $(1,1)/\sqrt2$ and $(1,-1)/\sqrt2$, for example, the vector $(3,1)$ has the coefficients $4/\sqrt2$ and $2/\sqrt2$, and indeed $(4/\sqrt2)(1,1)/\sqrt2+(2/\sqrt2)(1,-1)/\sqrt2=(2,2)+(1,-1)=(3,1)$. With an arbitrary family, finding the coefficients requires solving a linear system; this difference is the reason why orthonormal families are preferred whenever they are available, and \cref{ch:inner-products} shows how to construct them.

\subsection{Projection onto a Line and Distance to a Hyperplane}

\Cref{prop:projection-line} showed that the point of the line spanned by $\mathbf u\ne\mathbf 0$ closest to $\mathbf x$ is the \term{orthogonal projection}
\[
\mathbf p=\frac{\langle\mathbf u,\mathbf x\rangle}{\|\mathbf u\|^2}\,\mathbf u ,
\]
and that the residual $\mathbf x-\mathbf p$ is orthogonal to $\mathbf u$; \cref{fig:projection} illustrates it. The decomposition $\mathbf x=\mathbf p+(\mathbf x-\mathbf p)$ splits $\mathbf x$ into a part along $\mathbf u$ and a part orthogonal to $\mathbf u$, and by Pythagoras $\|\mathbf x\|^2=\|\mathbf p\|^2+\|\mathbf x-\mathbf p\|^2$. \Cref{sec:matrices-linear-maps} writes the map $\mathbf x\mapsto\mathbf p$ as a matrix, and \cref{ch:inner-products} replaces the line by a subspace of any dimension.

Orthogonality to a fixed vector is expressed by an equation. For $\mathbf a\ne\mathbf 0$ and a real number $b$, the set
\[
H\coloneqq\{\mathbf x\in\mathbb R^n: \langle\mathbf a,\mathbf x\rangle=b\}
\]
is a \term{hyperplane}, a line in the plane and a plane in space, and $\mathbf a$ is a \term{normal vector} of it: for two points $\mathbf x,\mathbf x'$ of $H$, $\langle\mathbf a,\mathbf x-\mathbf x'\rangle=b-b=0$, so every segment in $H$ is orthogonal to $\mathbf a$. The plane $x+y+z=1$ of \cref{sec:vectors-in-rn} has the normal vector $(1,1,1)$.

\begin{proposition}[Distance to a hyperplane]
\label{prop:hyperplane-distance}
Let $H=\{\mathbf x: \langle\mathbf a,\mathbf x\rangle=b\}$ with $\mathbf a\ne\mathbf 0$, and let $\mathbf y\in\mathbb R^n$. The point
\[
\mathbf z\coloneqq\mathbf y-\frac{\langle\mathbf a,\mathbf y\rangle-b}{\|\mathbf a\|^2}\,\mathbf a
\]
is the unique point of $H$ closest to $\mathbf y$, and the distance from $\mathbf y$ to $H$ is
\[
\|\mathbf y-\mathbf z\|=\frac{|\langle\mathbf a,\mathbf y\rangle-b|}{\|\mathbf a\|}.
\]
\end{proposition}

\begin{proof}
First, $\langle\mathbf a,\mathbf z\rangle=\langle\mathbf a,\mathbf y\rangle-(\langle\mathbf a,\mathbf y\rangle-b)=b$, so $\mathbf z\in H$. For every $\mathbf x\in H$, write $\mathbf y-\mathbf x=(\mathbf y-\mathbf z)+(\mathbf z-\mathbf x)$. The first part is a multiple of $\mathbf a$, and the second is orthogonal to $\mathbf a$, because $\langle\mathbf a,\mathbf z-\mathbf x\rangle=b-b=0$. By Pythagoras, $\|\mathbf y-\mathbf x\|^2=\|\mathbf y-\mathbf z\|^2+\|\mathbf z-\mathbf x\|^2$, which is smallest exactly for $\mathbf x=\mathbf z$. Finally, $\|\mathbf y-\mathbf z\|=|\langle\mathbf a,\mathbf y\rangle-b|\,\|\mathbf a\|/\|\mathbf a\|^2$.
\end{proof}

The distance of the origin from the plane $x+y+z=1$ is $1/\sqrt3$, attained at the point $(1,1,1)/3$, the centre of the triangle with the vertices $\mathbf e_1,\mathbf e_2,\mathbf e_3$. For a unit normal vector, $\|\mathbf a\|=1$, the number $\langle\mathbf a,\mathbf y\rangle-b$ is the \term{signed distance} of $\mathbf y$ from $H$: positive on the side into which $\mathbf a$ points, negative on the other side. Classifying points by the sign of an affine function $\langle\mathbf a,\mathbf y\rangle-b$ is the idea behind the linear classifiers of statistics and machine learning, and the orientation tests of computational geometry in \cref{sec:vector-numerics} are of the same kind.

\subsection{Other Norms}

The Euclidean norm measures length as a ruler does, but other ways of measuring the size of a vector are often more appropriate. When the components of $\mathbf x$ are errors, one may care about the largest of them; when they are amounts of money, about their total.

\begin{definition}[Norm]
\label{def:norm}
A \term{norm} on $\mathbb R^n$ is a function $\mathbf x\mapsto\|\mathbf x\|$ from $\mathbb R^n$ to $\mathbb R$ such that, for all $\mathbf x,\mathbf y\in\mathbb R^n$ and $t\in\mathbb R$,
\begin{enumerate}
\item $\|\mathbf x\|\ge0$, with equality only for $\mathbf x=\mathbf 0$;
\item $\|t\mathbf x\|=|t|\,\|\mathbf x\|$;
\item $\|\mathbf x+\mathbf y\|\le\|\mathbf x\|+\|\mathbf y\|$.
\end{enumerate}
The same definition applies to the matrices $\mathbb R^{m\times n}$ and to any other set with the arithmetic of \cref{prop:vector-rules}. The $1$-norm and the $\infty$-norm, or maximum norm, of $\mathbf x\in\mathbb R^n$ are
\[
\|\mathbf x\|_1\coloneqq\sum_{i=1}^n|x_i|,\qquad\|\mathbf x\|_\infty\coloneqq\max_{1\le i\le n}|x_i| ,
\]
and the Euclidean norm is also written $\|\mathbf x\|_2$ when it has to be told apart from them.
\end{definition}

\begin{proposition}[The three standard norms]
\label{prop:standard-norms}
$\|\cdot\|_1$ and $\|\cdot\|_\infty$ are norms on $\mathbb R^n$, and for every $\mathbf x\in\mathbb R^n$,
\[
\|\mathbf x\|_\infty\le\|\mathbf x\|_2\le\|\mathbf x\|_1\le\sqrt n\,\|\mathbf x\|_2\le n\|\mathbf x\|_\infty .
\]
Each inequality is an equality for some $\mathbf x\ne\mathbf 0$.
\end{proposition}

\begin{proof}
Positivity and homogeneity are clear for both. The triangle inequality follows from that for absolute values: $\sum_i|x_i+y_i|\le\sum_i|x_i|+\sum_i|y_i|$, and $|x_i+y_i|\le|x_i|+|y_i|\le\|\mathbf x\|_\infty+\|\mathbf y\|_\infty$ for every $i$, so that the maximum of the left side is also at most $\|\mathbf x\|_\infty+\|\mathbf y\|_\infty$. The first three inequalities of the chain are \cref{prop:norm}\,(v), and the last is $\|\mathbf x\|_2^2=\sum_ix_i^2\le n\|\mathbf x\|_\infty^2$. The vector $\mathbf e_1$ gives equality in the first two inequalities, and $(1,1,\ldots,1)$ in the last two.
\end{proof}

For $\mathbf x=(3,-4,0,1)$ the three norms are $\|\mathbf x\|_\infty=4$, $\|\mathbf x\|_2=\sqrt{26}=5.0990\ldots$ and $\|\mathbf x\|_1=8$. In the plane, $\|\mathbf x-\mathbf y\|_1$ is the distance travelled between $\mathbf x$ and $\mathbf y$ along a grid of streets, which is why it is also called the taxicab or Manhattan distance, and $\|\mathbf x-\mathbf y\|_\infty$ is the largest difference of a coordinate. \Cref{fig:unit-balls} shows the three \term{unit balls} $\{\mathbf x: \|\mathbf x\|\le1\}$: a diamond, a disc and a square, nested as the chain of inequalities predicts. The chain also shows that the three norms are interchangeable for questions of convergence: a sequence of vectors tends to $\mathbf 0$ in one of them if and only if it does in the others, since each is bounded by a constant multiple of each other. This is true of any two norms on $\mathbb R^n$, as Volume~II proves; the constants, however, depend on $n$, and in high dimensions they matter.

\begin{figure}
\sdfloat{%
\begin{tikzpicture}[scale=1.45, font=\scriptsize]
\fill[black!6] (-1,-1) rectangle (1,1);
\draw[black!45, line width=0.7pt] (-1,-1) rectangle (1,1);
\draw[black, dashed, line width=0.7pt] (0,0) circle (1);
\draw[black, line width=0.9pt] (1,0) -- (0,1) -- (-1,0) -- (0,-1) -- cycle;
\draw[line width=0.4pt, -{Latex[length=1.6mm]}] (-1.35,0) -- (1.45,0) node[right] {$x_1$};
\draw[line width=0.4pt, -{Latex[length=1.6mm]}] (0,-1.35) -- (0,1.4) node[left] {$x_2$};
\node[anchor=north west, inner sep=1pt] at (1.02,-0.02) {$1$};
\node[anchor=south east, inner sep=1pt] at (-0.02,1.02) {$1$};
\node[anchor=west] at (1.03,0.85) {$\|\mathbf x\|_\infty=1$};
\draw[line width=0.3pt] (0.893,0.45) -- (1.08,0.45);
\node[anchor=west] at (1.06,0.45) {$\|\mathbf x\|_2=1$};
\draw[line width=0.3pt] (0.65,-0.35) -- (1.08,-0.35);
\node[anchor=west] at (1.06,-0.35) {$\|\mathbf x\|_1=1$};
\end{tikzpicture}}{The unit balls of the $1$-norm (the diamond, solid), the Euclidean norm (the disc, dashed) and the $\infty$-norm (the square, grey and shaded) in the plane. The diamond lies inside the disc and the disc inside the square, because $\|\mathbf x\|_\infty\le\|\mathbf x\|_2\le\|\mathbf x\|_1$. All three are convex, as the unit ball of every norm is, by the triangle inequality.\label{fig:unit-balls}}
\end{figure}

\begin{note}{The \texorpdfstring{$p$}{p}-norms}
The three norms belong to a family. For $1\le p<\infty$ the \emph{$p$-norm} is $\|\mathbf x\|_p\coloneqq\bigl(\sum_i|x_i|^p\bigr)^{1/p}$; it gives $\|\cdot\|_1$ for $p=1$ and the Euclidean norm for $p=2$, and $\|\mathbf x\|_p\to\|\mathbf x\|_\infty$ as $p\to\infty$, since $\|\mathbf x\|_\infty\le\|\mathbf x\|_p\le n^{1/p}\|\mathbf x\|_\infty$. The triangle inequality for $\|\cdot\|_p$ is \emph{Minkowski's inequality}. It is derived from \emph{H\"older's inequality} $|\langle\mathbf x,\mathbf y\rangle|\le\|\mathbf x\|_p\|\mathbf y\|_q$ with $1/p+1/q=1$, which in turn follows from Young's inequality (\cref{cor:young}) and generalises the Cauchy--Schwarz inequality ($p=q=2$). As $p$ increases, the unit ball of $\|\cdot\|_p$ grows from the diamond to the square of \cref{fig:unit-balls}. For $0<p<1$ the same formula does not define a norm: the unit ball is not convex, and the triangle inequality fails, as $\mathbf e_1$ and $\mathbf e_2$ show. Hardy, Littlewood and P\'olya~\cite{hardy1952inequalities} and Horn and Johnson~\cite[Chap.~5]{horn2013} treat these inequalities and norms in detail; Volume~II uses them for the $L^p$ spaces of functions.
\end{note}

\section{Matrices and Linear Maps}
\label{sec:matrices-linear-maps}

\subsection{Matrices and the Matrix--Vector Product}

A matrix $A=(a_{ij})\in\mathbb R^{m\times n}$ is a rectangular array of $m$ rows and $n$ columns (\cref{def:matrix}). We write its columns as $\mathbf a_1,\ldots,\mathbf a_n\in\mathbb R^m$ and its rows as $\mathbf r_1^{\mathsf T},\ldots,\mathbf r_m^{\mathsf T}$, with $\mathbf r_i\in\mathbb R^n$, and display the two decompositions as
\[
A=\begin{bmatrix}\mathbf a_1&\mathbf a_2&\cdots&\mathbf a_n\end{bmatrix}=\begin{bmatrix}\mathbf r_1^{\mathsf T}\\\vdots\\\mathbf r_m^{\mathsf T}\end{bmatrix}.
\]
The \term{zero matrix}, all of whose entries are zero, is written $0$, and the \term{identity matrix} $I=I_n$ has the columns $\mathbf e_1,\ldots,\mathbf e_n$, so that its entries are $\delta_{ij}$. An $n\times1$ matrix is a column vector, a $1\times n$ matrix a row vector, and a $1\times1$ matrix is identified with its single entry. Matrices of the same shape are added and multiplied by scalars entry by entry, and these operations satisfy the eight rules of \cref{prop:vector-rules}, since $\mathbb R^{m\times n}$ is $\mathbb R^{mn}$ with the entries arranged in a rectangle.

\begin{definition}[Matrix--vector product]
\label{def:matrix-vector}
The product of $A\in\mathbb R^{m\times n}$ with columns $\mathbf a_1,\ldots,\mathbf a_n$ and of $\mathbf x\in\mathbb R^n$ is the linear combination of the columns with the components of $\mathbf x$ as coefficients,
\[
A\mathbf x\coloneqq x_1\mathbf a_1+x_2\mathbf a_2+\cdots+x_n\mathbf a_n\in\mathbb R^m .
\]
\end{definition}

The $i$th component of $x_1\mathbf a_1+\cdots+x_n\mathbf a_n$ is $\sum_jx_ja_{ij}$, so the definition agrees with the formula $(A\mathbf x)_i=\sum_ja_{ij}x_j$ of \cref{def:matrix}; the $i$th component is also the inner product $\langle\mathbf r_i,\mathbf x\rangle=\mathbf r_i^{\mathsf T}\mathbf x$ of the $i$th row with $\mathbf x$. The two readings, by columns and by rows, are worth keeping apart. The column reading answers the question ``what can $A\mathbf x$ be?'': the set of all products $A\mathbf x$ is the span of the columns, so the system $A\mathbf x=\mathbf b$ has a solution exactly when $\mathbf b$ is a combination of the columns. The row reading answers ``what does $A\mathbf x=\mathbf b$ say about $\mathbf x$?'': it says $\langle\mathbf r_i,\mathbf x\rangle=b_i$ for every $i$, so the solutions form the intersection of $m$ hyperplanes, one for each row. For the matrix with the rows $(1,1)$ and $(1,-1)$ and $\mathbf b=(3,1)$, the row reading asks for the intersection of the lines $x+y=3$ and $x-y=1$, the point $(2,1)$, and the column reading asks for the combination $2(1,1)+1(1,-1)=(3,1)$ of the columns. In particular $A\mathbf e_j=\mathbf a_j$: multiplying by a standard basis vector picks out a column.

\subsection{Linear Maps}

\Cref{def:linear-map} called a function $L\colon\mathbb R^n\to\mathbb R^m$ linear if $L(\mathbf x+\mathbf y)=L(\mathbf x)+L(\mathbf y)$ and $L(t\mathbf x)=tL(\mathbf x)$. A linear map preserves linear combinations, $L\bigl(\sum_jt_j\mathbf x_j\bigr)=\sum_jt_jL(\mathbf x_j)$, and therefore lines and planes: it maps the line $\{\mathbf p+t\mathbf v\}$ to the set $\{L(\mathbf p)+tL(\mathbf v)\}$, a line if $L(\mathbf v)\ne\mathbf 0$ and a single point otherwise. The following theorem, which sharpens \cref{prop:linear-matrix}, is the foundation of everything in this chapter.

\begin{theorem}[Matrices are linear maps]
\label{thm:matrices-linear-maps}
For $A\in\mathbb R^{m\times n}$ let $L_A\colon\mathbb R^n\to\mathbb R^m$ be the map $L_A(\mathbf x)\coloneqq A\mathbf x$.
\begin{enumerate}
\item $L_A$ is linear, and the $j$th column of $A$ is $L_A(\mathbf e_j)$.
\item For every linear map $L\colon\mathbb R^n\to\mathbb R^m$ there is exactly one matrix $A$ with $L=L_A$; its columns are $L(\mathbf e_1),\ldots,L(\mathbf e_n)$.
\item $L_{A+B}=L_A+L_B$ and $L_{tA}=tL_A$ for $A,B\in\mathbb R^{m\times n}$ and $t\in\mathbb R$, where the sum and multiples of maps are taken pointwise.
\end{enumerate}
So $A\mapsto L_A$ is a bijection from $\mathbb R^{m\times n}$ onto the set of linear maps from $\mathbb R^n$ to $\mathbb R^m$, compatible with addition and multiplication by scalars.
\end{theorem}

\begin{proof}
(i) By the rules of \cref{prop:vector-rules},
\[
A(\mathbf x+\mathbf y)=\sum_j(x_j+y_j)\mathbf a_j=\sum_jx_j\mathbf a_j+\sum_jy_j\mathbf a_j=A\mathbf x+A\mathbf y
\]
and $A(t\mathbf x)=\sum_j(tx_j)\mathbf a_j=t\,A\mathbf x$, and $A\mathbf e_j=\mathbf a_j$ because $\mathbf e_j$ has the single nonzero component $1$ in position $j$. (ii) Let $A$ have the columns $L(\mathbf e_j)$. For $\mathbf x=\sum_jx_j\mathbf e_j$, linearity gives $L(\mathbf x)=\sum_jx_jL(\mathbf e_j)=A\mathbf x$, so $L=L_A$. If also $L=L_B$, then the $j$th columns of $A$ and $B$ both equal $L(\mathbf e_j)$, by (i), so $A=B$. (iii) $(A+B)\mathbf x=\sum_jx_j(\mathbf a_j+\mathbf b_j)=A\mathbf x+B\mathbf x$ and $(tA)\mathbf x=\sum_jx_j(t\mathbf a_j)=t(A\mathbf x)$.
\end{proof}

The theorem gives a recipe that is used constantly: \emph{to find the matrix of a linear map, compute the images of the standard basis vectors and write them as columns}. A linear map is determined by what it does to $\mathbf e_1,\ldots,\mathbf e_n$, so a matrix is nothing but a table of these $n$ images. We shall freely identify a matrix with its linear map and speak of the matrix $A$ mapping $\mathbf x$ to $A\mathbf x$.

\subsection{Geometric Examples}

In the plane, the recipe produces the matrices of the familiar geometric maps, shown in \cref{fig:linear-maps}. The \term{rotation} by the angle $\varphi$ about the origin maps $\mathbf e_1$ to $(\cos\varphi,\sin\varphi)$ and $\mathbf e_2$ to $(-\sin\varphi,\cos\varphi)$, so its matrix is $R_\varphi$ of \cref{sec:vectors-rn}:
\[
R_\varphi=\begin{bmatrix}\cos\varphi&-\sin\varphi\\\sin\varphi&\cos\varphi\end{bmatrix}.
\]
That rotation is linear at all is a geometric fact: rotating the parallelogram spanned by $\mathbf x$ and $\mathbf y$ gives the parallelogram spanned by the rotated vectors, so the rotation of $\mathbf x+\mathbf y$ is the sum of the rotations. \Cref{sec:vectors-rn} confirmed it algebraically with the addition formulas.

\begin{figure}
\newcommand\sdFshape{(0,0) -- (0.25,0) -- (0.25,0.45) -- (0.6,0.45) -- (0.6,0.65) -- (0.25,0.65) -- (0.25,0.8) -- (0.75,0.8) -- (0.75,1) -- (0,1) -- cycle}
\newcommand\sdmappanel[2]{\begin{tikzpicture}[scale=1.05, font=\scriptsize]
#2
\draw[black!40, line width=0.4pt] \sdFshape;
\begin{scope}[cm={#1,(0,0)}]
\draw[black, dashed, line width=0.4pt] (0,0) rectangle (1,1);
\filldraw[fill=black!22, draw=black, line width=0.6pt] \sdFshape;
\end{scope}
\draw[line width=0.3pt, -{Latex[length=1.3mm]}] (-1,0) -- (2.05,0);
\draw[line width=0.3pt, -{Latex[length=1.3mm]}] (0,-0.45) -- (0,2.25);
\end{tikzpicture}}
\sdpanel{A}{\sdmappanel{0.7071,0.7071,-0.7071,0.7071}{}}\hfill
\sdpanel{B}{\sdmappanel{-0.5,0.8660,0.8660,0.5}{\draw[black!50, dotted, line width=0.6pt] (-0.25,-0.4330) -- (1.2,2.0785);}}\hfill
\sdpanel{C}{\sdmappanel{0.8,0.4,0.4,0.2}{\draw[black!50, dotted, line width=0.6pt] (-0.9,-0.45) -- (2,1);}}
\sdpanelrow
\sdpanel{D}{\sdmappanel{1,0,1,1}{}}\hfill
\sdpanel{E}{\sdmappanel{1.6,0,0,0.6}{}}\hfill
\sdpanel{F}{\sdmappanel{0.7071,0.7071,0,1.4142}{}}
\caption{Linear maps of the plane, applied to the letter F in the unit square (grey outline). The image of the letter is shaded, and the dashed parallelogram is the image of the unit square, spanned by the two columns of the matrix. (A) The rotation $R_{\pi/4}$. (B) The reflection $S_{\pi/3}$ in the dotted line at the angle $60$ degrees; the image is the mirror image of the letter. (C) The orthogonal projection onto the dotted line in the direction $(2,1)$; the whole letter is flattened onto a segment, and the square onto a segment of the same line. (D) The shear $K$ with $k=1$. (E) The scaling $\operatorname{diag}(1.6,0.6)$. (F) The shear of (D) followed by the rotation of (A), with the matrix $R_{\pi/4}K$ of \cref{sec:matrix-multiplication}.}
\label{fig:linear-maps}
\end{figure}

The \term{reflection} in the line through the origin at the angle $\alpha$ to the $x_1$-axis maps a unit vector at the angle $\beta$ to the unit vector at the angle $2\alpha-\beta$. Thus $\mathbf e_1$ (angle $0$) goes to $(\cos2\alpha,\sin2\alpha)$ and $\mathbf e_2$ (angle $\pi/2$) to $(\cos(2\alpha-\pi/2),\sin(2\alpha-\pi/2))=(\sin2\alpha,-\cos2\alpha)$, and the matrix is
\[
S_\alpha=\begin{bmatrix}\cos2\alpha&\sin2\alpha\\\sin2\alpha&-\cos2\alpha\end{bmatrix}.
\]
For $\alpha=0$ this is $\operatorname{diag}(1,-1)$, the reflection in the $x_1$-axis, which changes the sign of the second coordinate, and for $\alpha=\pi/4$ it exchanges the two coordinates. The \term{orthogonal projection} onto the same line, whose unit direction vector is $\mathbf u=(\cos\alpha,\sin\alpha)$, maps $\mathbf x$ to $\langle\mathbf u,\mathbf x\rangle\mathbf u$, so $\mathbf e_1\mapsto\cos\alpha\,\mathbf u$ and $\mathbf e_2\mapsto\sin\alpha\,\mathbf u$, and
\[
P_\alpha=\begin{bmatrix}\cos^2\alpha&\cos\alpha\sin\alpha\\\cos\alpha\sin\alpha&\sin^2\alpha\end{bmatrix}=\mathbf u\mathbf u^{\mathsf T}.
\]
By the double angle formulas, $S_\alpha=2P_\alpha-I$: the reflection of $\mathbf x$ is obtained by going from $\mathbf x$ to its projection $\mathbf p$ and then the same distance further, $\mathbf x+2(\mathbf p-\mathbf x)=2\mathbf p-\mathbf x$. Finally, the \term{shear} with the factor $k$ and the \term{scaling} with the factors $s_1$, $s_2$ have the matrices
\[
K=\begin{bmatrix}1&k\\0&1\end{bmatrix},\qquad\operatorname{diag}(s_1,s_2)=\begin{bmatrix}s_1&0\\0&s_2\end{bmatrix}.
\]
The shear moves each point horizontally by $k$ times its height, $K(x_1,x_2)=(x_1+kx_2,x_2)$, as a deck of cards is sheared by pushing its top sideways; the scaling stretches the two axes by different factors.

The figure makes visible what linearity means geometrically. Straight lines go to straight lines (or to points), parallel lines to parallel lines, the origin stays fixed, and a grid of congruent squares goes to a grid of congruent parallelograms, so the whole picture is determined by the image of one square, that is, by the two columns of the matrix. The projection is different from the other maps in one respect: it flattens the plane onto a line, distinct points such as $\mathbf x$ and $\mathbf x+(-1,2)$ go to the same point, and the map cannot be undone. \Cref{sec:inverse} returns to this.

Not every map that looks simple is linear. A \term{translation} $\mathbf x\mapsto\mathbf x+\mathbf b$ with $\mathbf b\ne\mathbf 0$ is not, since it moves the origin, and neither is $\mathbf x\mapsto\|\mathbf x\|$, since $\|-\mathbf x\|\ne-\|\mathbf x\|$ for $\mathbf x\ne\mathbf 0$. Maps of the form $\mathbf x\mapsto A\mathbf x+\mathbf b$ are called \term{affine}; the parametrisations of lines and planes in \cref{sec:vectors-in-rn} are examples. Computer graphics handles them with matrices all the same, by appending a component $1$ to every point of the plane, so that $\mathbf x\mapsto A\mathbf x+\mathbf b$ becomes the linear map of $\mathbb R^3$ with the matrix
\[
\begin{bmatrix}A&\mathbf b\\\mathbf 0^{\mathsf T}&1\end{bmatrix},\qquad\begin{bmatrix}A&\mathbf b\\\mathbf 0^{\mathsf T}&1\end{bmatrix}\begin{bmatrix}\mathbf x\\1\end{bmatrix}=\begin{bmatrix}A\mathbf x+\mathbf b\\1\end{bmatrix},
\]
a matrix written in blocks, as \cref{sec:block-matrices} explains; these \term{homogeneous coordinates} reappear in the geometry and graphics of Volume~XVI.

\subsection{Linear Maps beyond Geometry}

Linear maps are not only geometric. Many operations on lists of numbers are linear, and their matrices are found by the same recipe. The \term{mean} $\bar x=(x_1+\cdots+x_n)/n$ is the linear map from $\mathbb R^n$ to $\mathbb R$ with the $1\times n$ matrix $(1/n,\ldots,1/n)$. The \term{differences} of consecutive components, $\mathbf x\mapsto(x_2-x_1,x_3-x_2,\ldots,x_n-x_{n-1})$, form a linear map from $\mathbb R^n$ to $\mathbb R^{n-1}$ whose matrix has $-1$ on the diagonal and $1$ just above it; for $n=4$,
\[
\Delta=\begin{bmatrix}-1&1&0&0\\0&-1&1&0\\0&0&-1&1\end{bmatrix}.
\]
Divided by a step $h$, it computes the forward difference quotients of \cref{ch:differentiation} from the values of a function at equally spaced points, and it is the first example of a \emph{banded} matrix, studied in \cref{sec:special-matrices}.

A polynomial $p(x)=c_0+c_1x+\cdots+c_dx^d$ of degree at most $d$ is described by its vector of coefficients $\mathbf c=(c_0,\ldots,c_d)\in\mathbb R^{d+1}$, and many operations on polynomials become matrices. Differentiation maps $p$ to $p'(x)=c_1+2c_2x+\cdots+dc_dx^{d-1}$, and since $(p+q)'=p'+q'$ and $(tp)'=tp'$ (\cref{thm:diff-algebra}), it is linear. The images of the basis polynomials $1,x,x^2,x^3$ for $d=3$ are $0$, $1$, $2x$ and $3x^2$, with the coefficient vectors $\mathbf 0$, $\mathbf e_1$, $2\mathbf e_2$ and $3\mathbf e_3$ in $\mathbb R^4$, and these are the columns of
\[
D=\begin{bmatrix}0&1&0&0\\0&0&2&0\\0&0&0&3\\0&0&0&0\end{bmatrix}.
\]
For $p(x)=1+2x-x^2+4x^3$, $D(1,2,-1,4)=(2,-2,12,0)$, the coefficients of $p'(x)=2-2x+12x^2$. Evaluation at given points $x_0,\ldots,x_d$ is also linear: the vector of values $(p(x_0),\ldots,p(x_d))$ is $V\mathbf c$, where the \term{Vandermonde matrix} $V$ has the rows $(1,x_i,x_i^2,\ldots,x_i^d)$. Finding the polynomial that takes given values at given points, the interpolation problem, therefore means solving the linear system $V\mathbf c=\mathbf y$, which \cref{ch:linear-systems} and Volume~XI take up. Integration, multiplication by a fixed polynomial and the change of variable $x\mapsto x+1$ are linear in the same way. In each case, the recipe of \cref{thm:matrices-linear-maps} turns an operation on functions into a table of numbers, and that is how such operations are carried out by a computer.

\section{Matrix Multiplication}
\label{sec:matrix-multiplication}

\subsection{The Product as Composition}

If a linear map $B$ is followed by a linear map $A$, the composition is again linear, and so it has a matrix. The rule for computing this matrix is the definition of the matrix product; we derive it now, rather than postulating it as \cref{def:matrix-product} did.

\begin{theorem}[Composition and the matrix product]
\label{thm:product-composition}
Let $A\in\mathbb R^{m\times n}$ and $B\in\mathbb R^{n\times p}$, with columns $\mathbf b_1,\ldots,\mathbf b_p$. The composition $L_A\circ L_B\colon\mathbb R^p\to\mathbb R^m$, $\mathbf x\mapsto A(B\mathbf x)$, is linear, and its matrix is the product $AB\in\mathbb R^{m\times p}$ of \cref{def:matrix-product}, with the entries and columns
\[
(AB)_{ik}=\sum_{j=1}^na_{ij}b_{jk},\qquad AB=\begin{bmatrix}A\mathbf b_1&A\mathbf b_2&\cdots&A\mathbf b_p\end{bmatrix}.
\]
That is, $(AB)\mathbf x=A(B\mathbf x)$ for every $\mathbf x\in\mathbb R^p$.
\end{theorem}

\begin{proof}
For $\mathbf x,\mathbf y\in\mathbb R^p$ and scalars $s,t$, linearity of $L_B$ and then of $L_A$ gives $A(B(s\mathbf x+t\mathbf y))=A(sB\mathbf x+tB\mathbf y)=sA(B\mathbf x)+tA(B\mathbf y)$, so the composition is linear. By \cref{thm:matrices-linear-maps}, its matrix has the columns $A(B\mathbf e_k)=A\mathbf b_k$, and the $i$th component of $A\mathbf b_k$ is $\sum_ja_{ij}b_{jk}$. This is the matrix $AB$ of \cref{def:matrix-product}, and since a matrix represents its linear map, $(AB)\mathbf x=A(B\mathbf x)$ for every $\mathbf x$.
\end{proof}

The product $AB$ is defined only when the number of columns of $A$ equals the number of rows of $B$, because the maps must fit together: $B$ maps into $\mathbb R^n$, and $A$ must start there. A helpful mnemonic is that the shapes cancel like fractions, $(m\times n)(n\times p)=m\times p$. The order of the factors is the order of the composition: in $AB$, the map $B$ is applied first, as in $g\circ f$, and this order is forced by writing the vector on the right.

\Cref{fig:linear-maps}F is the shear $K$ followed by the rotation $R_{\pi/4}$. Its matrix is the product
\[
R_{\pi/4}K=\frac1{\sqrt2}\begin{bmatrix}1&-1\\1&1\end{bmatrix}\begin{bmatrix}1&1\\0&1\end{bmatrix}=\frac1{\sqrt2}\begin{bmatrix}1&0\\1&2\end{bmatrix},
\]
whose first column is $R_{\pi/4}\mathbf e_1$, since the shear fixes $\mathbf e_1$, and whose second column is the rotation of $K\mathbf e_2=(1,1)$, namely $(0,\sqrt2)$. The product of two rotations is a rotation, $R_\alpha R_\beta=R_{\alpha+\beta}$: multiplying out gives the entries $\cos\alpha\cos\beta-\sin\alpha\sin\beta=\cos(\alpha+\beta)$ and so on, and the addition formulas of \cref{thm:addition} are thereby exactly the statement that rotating by $\beta$ and then by $\alpha$ is rotating by $\alpha+\beta$.

\begin{theorem}[Rules of matrix multiplication]
\label{thm:product-rules}
Whenever the shapes fit,
\begin{enumerate}
\item $(AB)C=A(BC)$;
\item $A(B+C)=AB+AC$ and $(A+B)C=AC+BC$;
\item $t(AB)=(tA)B=A(tB)$ for $t\in\mathbb R$;
\item $I_mA=A=AI_n$ for $A\in\mathbb R^{m\times n}$, and $0A=0$, $A0=0$.
\end{enumerate}
\end{theorem}

\begin{proof}
(i) By \cref{thm:product-composition}, $(AB)C$ is the matrix of $L_{AB}\circ L_C=(L_A\circ L_B)\circ L_C$ and $A(BC)$ that of $L_A\circ(L_B\circ L_C)$. Composition of functions is associative (\cref{prop:composition}), so the two maps are equal, and by \cref{thm:matrices-linear-maps}\,(ii) so are their matrices. The same identity can be checked with entries: both sides have the $(i,l)$ entry $\sum_j\sum_ka_{ij}b_{jk}c_{kl}$, summed in two different orders. (ii) The $(i,k)$ entry of $A(B+C)$ is $\sum_ja_{ij}(b_{jk}+c_{jk})=\sum_ja_{ij}b_{jk}+\sum_ja_{ij}c_{jk}$, and similarly for the other law. (iii) Each of the three matrices has the entries $\sum_jta_{ij}b_{jk}$. (iv) $(I_mA)_{ik}=\sum_j\delta_{ij}a_{jk}=a_{ik}$, and similarly on the other side; the products with the zero matrix have only zero terms.
\end{proof}

Associativity is the most important of these rules. It allows us to write $ABC$ without parentheses, and, as \cref{sec:computing-matrices} shows, it gives the freedom to choose the cheaper of the two groupings. It is also the reason why the inconvenient-looking definition of the product is the right one: any definition under which products correspond to compositions inherits associativity from the composition of functions.

\subsection{The Product Is Not Commutative}

The rules of \cref{thm:product-rules} are those of the arithmetic of numbers, with one omission: there is no rule $AB=BA$. Even when both products are defined and have the same shape, they usually differ. With the rotation $R\coloneqq R_{\pi/2}$ by a right angle and the reflection $S\coloneqq\operatorname{diag}(1,-1)$ in the $x_1$-axis,
\[
RS=\begin{bmatrix}0&-1\\1&0\end{bmatrix}\begin{bmatrix}1&0\\0&-1\end{bmatrix}=\begin{bmatrix}0&1\\1&0\end{bmatrix},\qquad SR=\begin{bmatrix}1&0\\0&-1\end{bmatrix}\begin{bmatrix}0&-1\\1&0\end{bmatrix}=\begin{bmatrix}0&-1\\-1&0\end{bmatrix}.
\]
Reflecting in the $x_1$-axis and then rotating by a right angle is the reflection in the line $x_2=x_1$, while doing the same in the other order is the reflection in the line $x_2=-x_1$. Putting on socks and then shoes is not the same as putting on shoes and then socks, and the same is true of most pairs of operations. Some pairs do commute: two rotations, two diagonal matrices, any matrix and a multiple $tI$ of the identity, and any two powers of the same matrix. The \term{commutator} $AB-BA$ measures the failure; it is a basic object of quantum mechanics and of the theory of Lie groups.

A second difference from the arithmetic of numbers is that a product can be zero although neither factor is:
\[
\begin{bmatrix}1&1\\1&1\end{bmatrix}\begin{bmatrix}1&-1\\-1&1\end{bmatrix}=\begin{bmatrix}0&0\\0&0\end{bmatrix},\qquad N\coloneqq\begin{bmatrix}0&1\\0&0\end{bmatrix},\quad N^2=0 .
\]
Nonzero matrices with a zero product are \term{zero divisors}. Geometrically, the first factor on the left squeezes the plane onto the line spanned by $(1,1)$, and the second factor squeezes it onto the line spanned by $(1,-1)$, which the first one maps to $\mathbf 0$. As a consequence, matrices cannot be cancelled: $AB=AC$ with $A\ne0$ does not imply $B=C$, since $A(B-C)=0$ is possible with $B\ne C$. The familiar algebraic identities must be revised accordingly. By the distributive law, $(A+B)^2=A^2+AB+BA+B^2$, which equals $A^2+2AB+B^2$ only if $AB=BA$, and $(A+B)(A-B)=A^2-AB+BA-B^2$. Whenever $A$ and $B$ commute, the identities of ordinary algebra hold for them again; in particular the binomial theorem, which \cref{thm:binomial} proved for any two elements of a commutative ring, holds by the same induction:
\[
(A+B)^k=\sum_{j=0}^k\binom kjA^jB^{k-j}\qquad\text{if }AB=BA .
\]

\subsection{Four Ways to Read a Product}

The product $AB$ can be organised in several ways. They give the same matrix, but they suggest different ways of thinking and different algorithms.

\begin{theorem}[Four readings of the product]
\label{thm:four-readings}
Let $A\in\mathbb R^{m\times n}$ have the rows $\mathbf r_1^{\mathsf T},\ldots,\mathbf r_m^{\mathsf T}$ and the columns $\mathbf a_1,\ldots,\mathbf a_n$, and let $B\in\mathbb R^{n\times p}$ have the rows $\mathbf s_1^{\mathsf T},\ldots,\mathbf s_n^{\mathsf T}$ and the columns $\mathbf b_1,\ldots,\mathbf b_p$. Then:
\begin{enumerate}
\item (entries) $(AB)_{ik}=\mathbf r_i^{\mathsf T}\mathbf b_k=\langle\mathbf r_i,\mathbf b_k\rangle$;
\item (columns) the $k$th column of $AB$ is $A\mathbf b_k$, a combination of the columns of $A$ with the entries of $\mathbf b_k$ as coefficients;
\item (rows) the $i$th row of $AB$ is $\mathbf r_i^{\mathsf T}B=\sum_ja_{ij}\mathbf s_j^{\mathsf T}$, a combination of the rows of $B$ with the entries of $\mathbf r_i$ as coefficients;
\item (outer products) $AB=\sum_{j=1}^n\mathbf a_j\mathbf s_j^{\mathsf T}$, a sum of $n$ matrices of the form column times row.
\end{enumerate}
\end{theorem}

\begin{proof}
All four are rearrangements of the formula $(AB)_{ik}=\sum_ja_{ij}b_{jk}$. (i) The $i$th row of $A$ has the entries $a_{ij}$ and the $k$th column of $B$ the entries $b_{jk}$, $j=1,\ldots,n$. (ii) is \cref{thm:product-composition}, and $A\mathbf b_k=\sum_jb_{jk}\mathbf a_j$ by \cref{def:matrix-vector}. (iii) The $k$th entry of $\sum_ja_{ij}\mathbf s_j^{\mathsf T}$ is $\sum_ja_{ij}b_{jk}$. (iv) The outer product $\mathbf a_j\mathbf s_j^{\mathsf T}$ is the $m\times p$ matrix with the entries $a_{ij}b_{jk}$, and summing over $j$ gives $(AB)_{ik}$.
\end{proof}

\Cref{fig:product-views} illustrates the four readings, and a small example shows them at work. For
\[
A=\begin{bmatrix}1&2&0\\-1&3&1\end{bmatrix},\qquad B=\begin{bmatrix}2&1\\0&-1\\1&4\end{bmatrix},\qquad AB=\begin{bmatrix}2&-1\\-1&0\end{bmatrix},
\]
the entry $(AB)_{21}=-1$ is the inner product of the second row $(-1,3,1)$ of $A$ with the first column $(2,0,1)$ of $B$; the first column $(2,-1)$ of $AB$ is $2\mathbf a_1+0\mathbf a_2+1\mathbf a_3$; the first row $(2,-1)$ of $AB$ is $1\cdot(2,1)+2\cdot(0,-1)+0\cdot(1,4)$; and the sum of outer products is
\[
\begin{bmatrix}1\\-1\end{bmatrix}\begin{bmatrix}2&1\end{bmatrix}+\begin{bmatrix}2\\3\end{bmatrix}\begin{bmatrix}0&-1\end{bmatrix}+\begin{bmatrix}0\\1\end{bmatrix}\begin{bmatrix}1&4\end{bmatrix}=\begin{bmatrix}2&1\\-2&-1\end{bmatrix}+\begin{bmatrix}0&-2\\0&-3\end{bmatrix}+\begin{bmatrix}0&0\\1&4\end{bmatrix}=\begin{bmatrix}2&-1\\-1&0\end{bmatrix}.
\]

\begin{figure}
\newcommand\sdprodframe[1]{\begin{tikzpicture}[scale=1.0, font=\scriptsize, line width=0.4pt]
#1
\draw (0,0) rectangle (0.8,1.2);
\node at (1.05,0.6) {$=$};
\draw (1.3,0) rectangle (2.3,1.2);
\draw (2.45,0.1) rectangle (3.25,1.1);
\node[below] at (0.4,0) {$AB$};
\node[below] at (1.8,0) {$A$};
\node[below] at (2.85,0.1) {$B$};
\end{tikzpicture}}
\centering
\sdpanel{A}{\sdprodframe{\fill[black!45] (0.45,0.7) rectangle (0.6,0.85);
  \fill[black!25] (1.3,0.7) rectangle (2.3,0.85); \fill[black!25] (2.9,0.1) rectangle (3.05,1.1);}}\hspace{2cm}
\sdpanel{B}{\sdprodframe{\fill[black!25] (0.45,0) rectangle (0.6,1.2);
  \fill[black!10] (1.3,0) rectangle (2.3,1.2); \fill[black!25] (2.9,0.1) rectangle (3.05,1.1);}}
\sdpanelrow
\sdpanel{C}{\sdprodframe{\fill[black!25] (0,0.7) rectangle (0.8,0.85);
  \fill[black!25] (1.3,0.7) rectangle (2.3,0.85); \fill[black!10] (2.45,0.1) rectangle (3.25,1.1);}}\hspace{2cm}
\sdpanel{D}{\sdprodframe{\fill[black!10] (0,0) rectangle (0.8,1.2);
  \fill[black!25] (1.75,0) rectangle (1.9,1.2); \fill[black!25] (2.45,0.5) rectangle (3.25,0.65);}}
\caption{Four readings of the product $AB$ of an $m\times n$ and an $n\times p$ matrix (\cref{thm:four-readings}). (A) An entry of $AB$ is the inner product of a row of $A$ and a column of $B$. (B) A column of $AB$ is $A$ times a column of $B$, a combination of all columns of $A$. (C) A row of $AB$ is a row of $A$ times $B$, a combination of all rows of $B$. (D) The outer product of the $j$th column of $A$ and the $j$th row of $B$ is a full $m\times p$ matrix, and $AB$ is the sum of these $n$ layers.}
\label{fig:product-views}
\end{figure}

Each reading is useful in its own place. The entry reading is the one usually taught first, and it is the right one for computing a single entry. The column reading says that every column of $AB$ lies in the span of the columns of $A$, which \cref{ch:vector-spaces} uses to compare ranks, and it is the natural way to multiply when the matrices are stored column by column. The row reading is the transposed statement. The outer product reading writes $AB$ as a sum of $n$ \term{rank-one matrices}, matrices of the form $\mathbf u\mathbf v^{\mathsf T}$, all of whose columns are multiples of the single vector $\mathbf u$. It is the reading behind Gaussian elimination (\cref{ch:linear-systems}), where each step subtracts a rank-one matrix, and behind the singular value decomposition (Chapter~21), which writes a matrix as a sum of rank-one matrices ordered by importance; it also shows that the product can be accumulated in $n$ steps, each of which reads one column of $A$ and one row of $B$.

\subsection{Powers and Polynomials of a Matrix}

A square matrix can be multiplied by itself. The \term{powers} of $A\in\mathbb R^{n\times n}$ are $A^0\coloneqq I$ and $A^{k+1}\coloneqq AA^k$; by associativity, $A^jA^k=A^{j+k}$ and $(A^j)^k=A^{jk}$, exactly as for numbers, and in particular any two powers of $A$ commute. The rule $(AB)^k=A^kB^k$, however, fails unless $A$ and $B$ commute, since $(AB)^2=ABAB$. For a polynomial $p(x)=c_0+c_1x+\cdots+c_dx^d$ we define
\[
p(A)\coloneqq c_0I+c_1A+c_2A^2+\cdots+c_dA^d .
\]
Because powers of $A$ commute, products of such matrices multiply out like polynomials: if $r=pq$ is the product of two polynomials, then $r(A)=p(A)q(A)=q(A)p(A)$. Indeed, by the distributive law, $p(A)q(A)=\sum_{j,k}c_jd_kA^{j+k}$, where $d_k$ are the coefficients of $q$, and collecting the terms with $j+k=l$ gives the coefficients of $r$. This is the start of the theory of matrix functions in Chapter~22, where the exponential $e^A=\sum_kA^k/k!$ is defined by the series.

Powers of a matrix describe repeated application of a linear map, and they appear wherever a linear rule is iterated. The rotations satisfy $R_\varphi^k=R_{k\varphi}$, by the product of rotations computed above. The nilpotent matrix $N$ of the previous subsection has $N^2=0$, so the binomial theorem, applied to the commuting matrices $I$ and $N$, gives $(I+N)^k=I+kN$: the $k$th power of a shear is the shear with $k$ times the factor. A recurrence also becomes a matrix power. The Fibonacci numbers of \cref{sec:gcd}, with $F_0=0$, $F_1=1$ and $F_{k+1}=F_k+F_{k-1}$, satisfy
\[
\begin{bmatrix}F_{k+1}\\F_k\end{bmatrix}=\begin{bmatrix}1&1\\1&0\end{bmatrix}\begin{bmatrix}F_k\\F_{k-1}\end{bmatrix},\qquad\text{so}\qquad\begin{bmatrix}1&1\\1&0\end{bmatrix}^k=\begin{bmatrix}F_{k+1}&F_k\\F_k&F_{k-1}\end{bmatrix}\quad(k\ge1),
\]
the second identity by induction on $k$, since multiplying the right side by the matrix gives the rows $(F_{k+1}+F_k,F_{k+1})=(F_{k+2},F_{k+1})$ and $(F_k+F_{k-1},F_k)=(F_{k+1},F_k)$. For $k=10$ the power is the matrix with the rows $(89,55)$ and $(55,34)$. Comparing the entries of $F^{j+k}=F^jF^k$, where $F$ is the matrix on the left, gives identities such as $F_{j+k}=F_{j+1}F_k+F_jF_{k-1}$, and the power $F^k$ can be computed with about $2\log_2k$ matrix products by repeated squaring, $F^{2k}=(F^k)^2$. Chapter~27 solves linear recurrences in general, and \cref{ch:eigenvalues} explains the growth of $F^k$, like $\phi^k$, through the eigenvalues of $F$.

\subsection{The Cost of a Product and the Order of the Loops}

Computing one entry of $AB$ by the entry reading takes $n$ multiplications and $n-1$ additions, so the product of an $m\times n$ and an $n\times p$ matrix takes $mnp$ multiplications and $mp(n-1)$ additions. For square matrices of order $n$ this is $n^3$ multiplications, about $2n^3$ arithmetic operations in all; doubling the size multiplies the work by eight. \Cref{alg:matmul} states the computation with three nested loops.

\begin{algorithm}
\begin{algorithmic}[1]
\Require $A=(a_{ij})\in\mathbb R^{m\times n}$ and $B=(b_{jk})\in\mathbb R^{n\times p}$
\Ensure $C=AB\in\mathbb R^{m\times p}$
\Name{MatMul}{A,B}
\State $C\gets0$ \Comment{the $m\times p$ zero matrix}
\For{$i=1,\ldots,m$}
  \For{$k=1,\ldots,p$}
    \For{$j=1,\ldots,n$}
      \State $c_{ik}\gets c_{ik}+a_{ij}b_{jk}$
    \EndFor
  \EndFor
\EndFor
\State \Return $C$
\end{algorithmic}
\caption{The matrix product by three nested loops, in the order $i,k,j$. The three \textbf{for} lines may be arranged in any of the six possible orders without changing the result (\cref{prop:loop-orders}). The cost is $mnp$ multiplications and as many additions.}
\label{alg:matmul}
\end{algorithm}

The three loops are independent of each other, and they can be nested in any of the $3!=6$ orders. The innermost loop then computes something different in each case. With $j$ innermost, as in the algorithm, it computes an inner product of a row of $A$ and a column of $B$, the entry reading. With $i$ innermost, it adds the multiple $b_{jk}\mathbf a_j$ of a column of $A$ to the $k$th column of $C$, the column reading; with $k$ innermost, it adds the multiple $a_{ij}\mathbf s_j^{\mathsf T}$ of a row of $B$ to the $i$th row of $C$, the row reading. With $j$ outermost, the two inner loops add the outer product $\mathbf a_j\mathbf s_j^{\mathsf T}$ to $C$, the outer product reading. The order does not change the result even in floating-point arithmetic, where the order of additions usually matters, and the reason is simple.

\begin{proposition}[Loop orders]
\label{prop:loop-orders}
In each of the six nestings of the loops of \cref{alg:matmul}, every entry $c_{ik}$ receives the same sequence of updates $c_{ik}\gets c_{ik}+a_{ij}b_{jk}$, in the order $j=1,2,\ldots,n$. Consequently all six orders produce the same matrix, also when every operation is rounded.
\end{proposition}

\begin{proof}
In a nesting of three loops, each running upwards, the triples $(i,j,k)$ are visited in lexicographic order with respect to the order in which the loop variables are nested: a triple comes before another if, in the first loop variable where they differ, it has the smaller value. Two triples $(i,j,k)$ and $(i,j',k)$ with the same $i$ and $k$ differ only in $j$, so the one with the smaller $j$ comes first, whatever the nesting. An update with indices $(i,j,k)$ changes only the entry $c_{ik}$, so the computation of $c_{ik}$ is the sequence of updates with indices $(i,1,k),(i,2,k),\ldots,(i,n,k)$ in this order, starting from $0$, independently of the nesting. In floating-point arithmetic each update is the same rounded operation applied to the same operands, so the computed values agree as well.
\end{proof}

The six orders differ in the order in which they read the entries of $A$ and $B$ from memory. A computer stores a matrix as one long list, either row by row or column by column, and it reads consecutive entries much faster than scattered ones, because the memory is organised in a hierarchy of caches that transfer blocks of neighbouring entries at a time. An order whose innermost loop runs along the stored rows or columns can be several times faster than one that jumps through memory, although the arithmetic is identical~\cite{dongarra1984}, \cite[Sect.~1.1]{golub2013}. High-performance libraries go further: they split the matrices into blocks that fit into the caches and multiply the blocks with carefully tuned kernels~\cite{goto2008}; \cref{sec:block-matrices} shows why computing with blocks is legitimate. Volume~IV explains the memory hierarchy, and Volume~VIII the blocked algorithms. The number $n^3$ itself is not a law of nature either: there are algorithms that multiply matrices with asymptotically fewer operations, which the note on Strassen's algorithm in \cref{sec:block-matrices} describes.

\section{The Transpose}
\label{sec:transpose}

\subsection{Rules of Transposition}

The transpose $A^{\mathsf T}\in\mathbb R^{n\times m}$ of $A\in\mathbb R^{m\times n}$ has the entries $(A^{\mathsf T})_{ji}=a_{ij}$ (\cref{def:matrix-product}): its rows are the columns of $A$, and its columns are the rows of $A$. For example,
\[
\begin{bmatrix}1&2&0\\-1&3&1\end{bmatrix}^{\mathsf T}=\begin{bmatrix}1&-1\\2&3\\0&1\end{bmatrix},\qquad\begin{bmatrix}x_1\\\vdots\\x_n\end{bmatrix}^{\mathsf T}=\begin{bmatrix}x_1&\cdots&x_n\end{bmatrix}.
\]
Geometrically, a square matrix is transposed by reflecting it in its \term{main diagonal}, the entries $a_{11},a_{22},\ldots,a_{nn}$, which stay in place. Transposition is a purely formal operation on the table of entries, but it has a meaning for the linear map, which part~(v) of the following proposition reveals.

\begin{proposition}[Rules of transposition]
\label{prop:transpose-rules}
Whenever the shapes fit,
\begin{enumerate}
\item $(A^{\mathsf T})^{\mathsf T}=A$, $(A+B)^{\mathsf T}=A^{\mathsf T}+B^{\mathsf T}$ and $(tA)^{\mathsf T}=tA^{\mathsf T}$;
\item $(AB)^{\mathsf T}=B^{\mathsf T}A^{\mathsf T}$, and more generally $(A_1A_2\cdots A_k)^{\mathsf T}=A_k^{\mathsf T}\cdots A_2^{\mathsf T}A_1^{\mathsf T}$;
\item $\mathbf x^{\mathsf T}A\mathbf y=\mathbf y^{\mathsf T}A^{\mathsf T}\mathbf x$ for $A\in\mathbb R^{m\times n}$, $\mathbf x\in\mathbb R^m$ and $\mathbf y\in\mathbb R^n$;
\item $(\mathbf u\mathbf v^{\mathsf T})^{\mathsf T}=\mathbf v\mathbf u^{\mathsf T}$;
\item $\langle A\mathbf x,\mathbf y\rangle=\langle\mathbf x,A^{\mathsf T}\mathbf y\rangle$ for all $\mathbf x\in\mathbb R^n$, $\mathbf y\in\mathbb R^m$, and $A^{\mathsf T}$ is the only matrix $B\in\mathbb R^{n\times m}$ with $\langle A\mathbf x,\mathbf y\rangle=\langle\mathbf x,B\mathbf y\rangle$ for all $\mathbf x$ and $\mathbf y$.
\end{enumerate}
\end{proposition}

\begin{proof}
(i) is immediate from the definition. (ii) The $(k,i)$ entry of $(AB)^{\mathsf T}$ is $(AB)_{ik}=\sum_ja_{ij}b_{jk}=\sum_j(B^{\mathsf T})_{kj}(A^{\mathsf T})_{ji}=(B^{\mathsf T}A^{\mathsf T})_{ki}$, as in \cref{prop:matrix-algebra}\,(iii). The general case follows by induction on $k$, writing $A_1\cdots A_k=(A_1\cdots A_{k-1})A_k$. (iii) The left side is a $1\times1$ matrix, which equals its own transpose, and by (ii) this transpose is $\mathbf y^{\mathsf T}A^{\mathsf T}\mathbf x$. (iv) is (ii) for the product of the $n\times1$ matrix $\mathbf u$ and the $1\times n$ matrix $\mathbf v^{\mathsf T}$. (v) By (iii), $\langle A\mathbf x,\mathbf y\rangle=\mathbf y^{\mathsf T}A\mathbf x=\mathbf x^{\mathsf T}A^{\mathsf T}\mathbf y=\langle\mathbf x,A^{\mathsf T}\mathbf y\rangle$. If $B$ has the same property, then choosing $\mathbf x=\mathbf e_j$ and $\mathbf y=\mathbf e_i$ gives $a_{ij}=\langle A\mathbf e_j,\mathbf e_i\rangle=\langle\mathbf e_j,B\mathbf e_i\rangle=b_{ji}$, so $B=A^{\mathsf T}$.
\end{proof}

Part (ii), the \term{reversal rule}, is the same rule as for inverse functions, $(g\circ f)^{-1}=f^{-1}\circ g^{-1}$ (\cref{prop:inverse-composition}): undoing or transposing a sequence of operations reverses their order, as taking off shoes and socks reverses the order of putting them on. Part (v) characterises the transpose without coordinates: $A^{\mathsf T}$ is the map that moves $A$ from one side of the inner product to the other, the \term{adjoint} of $A$. This is how the transpose is generalised to the inner product spaces of \cref{ch:inner-products} and to the integral operators of later volumes, and it is also how it arises in the reverse mode of automatic differentiation (\cref{sec:chain-rule-n}). There the gradient of a composition is computed by the chain rule as $\nabla(f\circ\mathbf g)(\mathbf a)=D\mathbf g(\mathbf a)^{\mathsf T}\nabla f(\mathbf g(\mathbf a))$, by \eqref{eq:chain-gradient} and the reversal rule: the transposed Jacobian matrices are applied in the reverse order of the computation.

\subsection{Symmetric and Antisymmetric Matrices}

A square matrix is \term{symmetric} if $A^{\mathsf T}=A$, that is, $a_{ij}=a_{ji}$ for all $i,j$, and \term{antisymmetric}, or skew-symmetric, if $A^{\mathsf T}=-A$. The diagonal entries of an antisymmetric matrix vanish, since $a_{ii}=-a_{ii}$. Hessian matrices are symmetric, by Schwarz's theorem (\cref{thm:schwarz}), as are the adjacency matrices of undirected graphs (\cref{sec:graphs}) and the distance tables of road atlases. The cross product with a fixed vector $\mathbf a\in\mathbb R^3$ is an antisymmetric linear map: from the formula for $\mathbf a\times\mathbf x$ in \cref{sec:surfaces},
\[
\mathbf a\times\mathbf x=\begin{bmatrix}0&-a_3&a_2\\a_3&0&-a_1\\-a_2&a_1&0\end{bmatrix}\mathbf x .
\]

\begin{proposition}[Symmetric and antisymmetric parts]
\label{prop:symmetric-parts}
Every $A\in\mathbb R^{n\times n}$ is in exactly one way the sum $A=A_{\mathrm s}+A_{\mathrm a}$ of a symmetric matrix $A_{\mathrm s}$ and an antisymmetric matrix $A_{\mathrm a}$, namely
\[
A_{\mathrm s}=\frac{A+A^{\mathsf T}}2,\qquad A_{\mathrm a}=\frac{A-A^{\mathsf T}}2 .
\]
Moreover $\langle\mathbf x,K\mathbf x\rangle=0$ for every antisymmetric $K$ and every $\mathbf x$, so that $\langle\mathbf x,A\mathbf x\rangle=\langle\mathbf x,A_{\mathrm s}\mathbf x\rangle$.
\end{proposition}

\begin{proof}
By \cref{prop:transpose-rules}\,(i), $A_{\mathrm s}^{\mathsf T}=(A^{\mathsf T}+A)/2=A_{\mathrm s}$ and $A_{\mathrm a}^{\mathsf T}=(A^{\mathsf T}-A)/2=-A_{\mathrm a}$, and their sum is $A$. If also $A=S+K$ with $S$ symmetric and $K$ antisymmetric, then $A^{\mathsf T}=S-K$, so $S=(A+A^{\mathsf T})/2$ and $K=(A-A^{\mathsf T})/2$. For antisymmetric $K$, the number $\langle\mathbf x,K\mathbf x\rangle=\mathbf x^{\mathsf T}K\mathbf x$ equals its transpose $\mathbf x^{\mathsf T}K^{\mathsf T}\mathbf x=-\mathbf x^{\mathsf T}K\mathbf x$, so it is zero.
\end{proof}

For example,
\[
\begin{bmatrix}1&2\\4&3\end{bmatrix}=\begin{bmatrix}1&3\\3&3\end{bmatrix}+\begin{bmatrix}0&-1\\1&0\end{bmatrix}.
\]
The antisymmetric part here is the rotation $R_{\pi/2}$, which maps every vector to an orthogonal one, in accordance with $\langle\mathbf x,K\mathbf x\rangle=0$. The proposition explains why \cref{def:definite} considered definiteness for symmetric matrices only: the quadratic form $\langle\mathbf x,A\mathbf x\rangle$ sees only the symmetric part of $A$, so nothing is lost by assuming $A$ symmetric from the start.

\subsection{Gram Matrices}

The products $A^{\mathsf T}A$ and $AA^{\mathsf T}$, which are defined for every matrix $A$, appear throughout the series, beginning with the normal equations $X^{\mathsf T}X\boldsymbol\beta=X^{\mathsf T}\mathbf y$ of least squares (\cref{sec:extrema-n}).

\begin{proposition}[Gram matrices]
\label{prop:gram}
Let $A\in\mathbb R^{m\times n}$ have the columns $\mathbf a_1,\ldots,\mathbf a_n$. Then $G\coloneqq A^{\mathsf T}A\in\mathbb R^{n\times n}$ has the entries $g_{jk}=\langle\mathbf a_j,\mathbf a_k\rangle$; it is symmetric, and
\[
\langle\mathbf x,A^{\mathsf T}A\mathbf x\rangle=\|A\mathbf x\|^2\ge0\qquad\text{for all }\mathbf x\in\mathbb R^n ,
\]
so that $A^{\mathsf T}A$ is positive semidefinite. It is positive definite if and only if $A\mathbf x=\mathbf 0$ only for $\mathbf x=\mathbf 0$. The same holds for $AA^{\mathsf T}$, with the rows of $A$ in place of the columns.
\end{proposition}

\begin{proof}
The $(j,k)$ entry of $A^{\mathsf T}A$ is the $j$th row of $A^{\mathsf T}$, which is $\mathbf a_j^{\mathsf T}$, times the $k$th column of $A$, by \cref{thm:four-readings}\,(i). Symmetry follows from $(A^{\mathsf T}A)^{\mathsf T}=A^{\mathsf T}(A^{\mathsf T})^{\mathsf T}=A^{\mathsf T}A$, and the identity from \cref{prop:transpose-rules}\,(v): $\langle\mathbf x,A^{\mathsf T}A\mathbf x\rangle=\langle A\mathbf x,A\mathbf x\rangle$. This is positive for every $\mathbf x\ne\mathbf 0$ exactly when $A\mathbf x\ne\mathbf 0$ for every $\mathbf x\ne\mathbf 0$. For $AA^{\mathsf T}$ apply the result to $A^{\mathsf T}$, whose columns are the rows of $A$.
\end{proof}

The matrix $A^{\mathsf T}A$ is the \term{Gram matrix} of the vectors $\mathbf a_1,\ldots,\mathbf a_n$, the table of all their inner products. It contains all the geometry of the family, the lengths and the angles, but forgets its position in space. For the columns $(1,1,0)$ and $(1,0,1)$,
\[
A=\begin{bmatrix}1&1\\1&0\\0&1\end{bmatrix},\qquad A^{\mathsf T}A=\begin{bmatrix}2&1\\1&2\end{bmatrix},\qquad AA^{\mathsf T}=\begin{bmatrix}2&1&1\\1&1&0\\1&0&1\end{bmatrix};
\]
the two vectors have length $\sqrt2$ and the angle $\pi/3$, since $\cos\theta=1/2$. The Gram matrix of an orthonormal family is the identity, and conversely: \emph{the columns of $Q$ are orthonormal if and only if $Q^{\mathsf T}Q=I$}, a criterion that \cref{sec:special-matrices} uses. In statistics, if the columns of $X$ hold centred measurements of several variables, $X^{\mathsf T}X/(N-1)$ is the sample covariance matrix (Volume~III), and its positive semidefiniteness reflects the fact that no combination of the variables can have a negative variance.

\section{The Inverse}
\label{sec:inverse}

\subsection{Left, Right and Two-Sided Inverses}

The inverse undoes a linear map. For numbers, $a\ne0$ has the inverse $a^{-1}$ with $a^{-1}a=aa^{-1}=1$; for matrices, the two products are different, and a matrix that is not square can have an inverse on one side only.

\begin{definition}[Inverses]
\label{def:matrix-inverse}
Let $A\in\mathbb R^{m\times n}$. A matrix $B\in\mathbb R^{n\times m}$ is a \term{left inverse} of $A$ if $BA=I_n$, and a \term{right inverse} if $AB=I_m$. A square matrix $A\in\mathbb R^{n\times n}$ is \term{invertible}, or \term{nonsingular}, if it has a two-sided inverse $B$, with $AB=BA=I$; this $B$ is the \term{inverse} $A^{-1}$. A square matrix that is not invertible is \term{singular}.
\end{definition}

\begin{proposition}[One-sided inverses]
\label{prop:one-sided-inverses}
Let $A\in\mathbb R^{m\times n}$.
\begin{enumerate}
\item If $A$ has a left inverse $B$ and a right inverse $C$, then $B=C$. In particular a two-sided inverse is unique, and if $A$ has a two-sided inverse, it has no other left or right inverse.
\item If $A$ has a left inverse, then $A\mathbf x=\mathbf 0$ only for $\mathbf x=\mathbf 0$, and $A\mathbf x=\mathbf b$ has at most one solution for every $\mathbf b$: the map $L_A$ is injective.
\item If $A$ has a right inverse $C$, then $A\mathbf x=\mathbf b$ has the solution $\mathbf x=C\mathbf b$ for every $\mathbf b$: the map $L_A$ is surjective.
\item A square matrix $A$ is invertible if and only if $L_A$ is bijective, and then $L_{A^{-1}}$ is the inverse function of $L_A$.
\end{enumerate}
\end{proposition}

\begin{proof}
(i) $B=BI_m=B(AC)=(BA)C=I_nC=C$, by associativity. (ii) If $BA=I$ and $A\mathbf x=A\mathbf y$, then $\mathbf x=BA\mathbf x=BA\mathbf y=\mathbf y$. (iii) $A(C\mathbf b)=(AC)\mathbf b=\mathbf b$. (iv) If $A$ is invertible, $L_A$ is injective and surjective by (ii) and (iii), and $L_{A^{-1}}\circ L_A=L_{A^{-1}A}=\operatorname{id}$, and likewise in the other order. Conversely, if $L_A$ is bijective, its inverse function $g$ is linear: for $\mathbf u=A\mathbf x$ and $\mathbf v=A\mathbf y$, $g(s\mathbf u+t\mathbf v)=g(A(s\mathbf x+t\mathbf y))=s\mathbf x+t\mathbf y=sg(\mathbf u)+tg(\mathbf v)$. By \cref{thm:matrices-linear-maps}, $g=L_B$ for a matrix $B$, and $L_{AB}=L_A\circ L_B=\operatorname{id}=L_{BA}$ gives $AB=BA=I$.
\end{proof}

The column vector $\mathbf a=(1,2)$, a $2\times1$ matrix, shows the asymmetry of the two sides. The map $L_{\mathbf a}\colon t\mapsto t\mathbf a$ from $\mathbb R$ to $\mathbb R^2$ is injective, and every row $(b_1,b_2)$ with $b_1+2b_2=1$ is a left inverse, for example $(1,0)$, $(0,1/2)$ and $(1,2)/5$; each of them recovers $t$ from the point $t\mathbf a$, but each in a different way, by reading off the first coordinate, half the second coordinate, or the projection onto the line. There is no right inverse: for every row $\mathbf b^{\mathsf T}=(b_1,b_2)$, the product $\mathbf a\mathbf b^{\mathsf T}$ has the rows $(b_1,b_2)$ and $(2b_1,2b_2)$, the second twice the first, and so it is not $I_2$. Correspondingly, $L_{\mathbf a}$ is not surjective: its image is a line. Chapter~21 singles out the left inverse $(1,2)/5$, the one given by orthogonal projection, as the \emph{pseudoinverse}.

For square matrices the asymmetry disappears. This is a genuinely deeper fact, because it depends on the dimension being the same on both sides.

\begin{theorem}[One-sided inverses of square matrices]
\label{thm:one-sided-square}
Let $A,B\in\mathbb R^{n\times n}$. If $BA=I$, then also $AB=I$, so that $A$ is invertible with $A^{-1}=B$. The same holds if $AB=I$.
\end{theorem}

\begin{proof}
If $BA=I$, then $A\mathbf x=\mathbf 0$ only for $\mathbf x=\mathbf 0$, by \cref{prop:one-sided-inverses}\,(ii). For a \emph{square} matrix with this property, the system $A\mathbf x=\mathbf b$ has a solution for every $\mathbf b$; this is part of \cref{thm:invertible-matrices}\,(i), and \cref{ch:linear-systems} proves it by Gaussian elimination, which reduces a square matrix without nonzero solutions of $A\mathbf x=\mathbf 0$ to the identity by row operations, each of which is undone by another. Solving $A\mathbf c_k=\mathbf e_k$ for $k=1,\ldots,n$ gives a matrix $C$ with the columns $\mathbf c_k$ and $AC=I$, and by \cref{prop:one-sided-inverses}\,(i), $B=C$. If $AB=I$, apply the first case to the pair $(B,A)$: $B$ has the left inverse $A$, so $BA=I$ as well. See also Horn and Johnson~\cite[Sect.~0.5]{horn2013} and Strang~\cite[Sect.~2.5]{strang2016}.
\end{proof}

The theorem is false in infinite dimensions, which shows that it cannot follow from the algebraic rules alone. On the space of real sequences, the right shift $S(x_0,x_1,\ldots)\coloneqq(0,x_0,x_1,\ldots)$ and the left shift $T(x_0,x_1,\ldots)\coloneqq(x_1,x_2,\ldots)$ are linear and satisfy $T\circ S=\operatorname{id}$, but $S\circ T$ replaces $x_0$ by $0$ and is not the identity. In finite dimensions, counting dimensions (\cref{ch:linear-maps}) shows that an injective linear map from $\mathbb R^n$ to itself must be surjective.

\subsection{Rules for Inverses}

\begin{proposition}[Rules for inverses]
\label{prop:inverse-rules}
Let $A,B\in\mathbb R^{n\times n}$ be invertible and $t\ne0$. Then $A^{-1}$, $AB$, $A^{\mathsf T}$ and $tA$ are invertible, and
\begin{enumerate}
\item $(A^{-1})^{-1}=A$;
\item $(AB)^{-1}=B^{-1}A^{-1}$;
\item $(A^{\mathsf T})^{-1}=(A^{-1})^{\mathsf T}$;
\item $(tA)^{-1}=t^{-1}A^{-1}$.
\end{enumerate}
With $A^{-k}\coloneqq(A^{-1})^k$ for $k\in\mathbb N$, the rules $A^jA^k=A^{j+k}$ and $(A^j)^k=A^{jk}$ hold for all integers $j,k$.
\end{proposition}

\begin{proof}
In each case it suffices to check that the proposed inverse is a two-sided inverse, by \cref{prop:one-sided-inverses}\,(i). (i) $AA^{-1}=A^{-1}A=I$ says that $A$ is the inverse of $A^{-1}$. (ii) $(AB)(B^{-1}A^{-1})=A(BB^{-1})A^{-1}=AA^{-1}=I$, and similarly $(B^{-1}A^{-1})(AB)=I$. (iii) By the reversal rule, $A^{\mathsf T}(A^{-1})^{\mathsf T}=(A^{-1}A)^{\mathsf T}=I^{\mathsf T}=I$, and similarly in the other order. (iv) $(tA)(t^{-1}A^{-1})=(tt^{-1})AA^{-1}=I$. For the last statement, (ii) applied repeatedly shows that $A^k$ is invertible with the inverse $(A^{-1})^k$, and the rules for nonnegative exponents extend to negative ones by cancelling factors $AA^{-1}=I$.
\end{proof}

The reversal rule (ii) appears again: to undo $B$ followed by $A$, first undo $A$, then undo $B$. In particular, the invertible $n\times n$ matrices are closed under products and inverses; with the identity as neutral element they form a \term{group} in the sense of algebra, the \term{general linear group}.

The geometric maps of \cref{sec:matrices-linear-maps} have evident inverses. The rotation by $\varphi$ is undone by the rotation by $-\varphi$, so $R_\varphi^{-1}=R_{-\varphi}$, and since $\cos(-\varphi)=\cos\varphi$ and $\sin(-\varphi)=-\sin\varphi$ this is also $R_\varphi^{\mathsf T}$. A reflection is its own inverse, $S_\alpha^2=I$. The shear by $k$ is undone by the shear by $-k$, and a scaling $\operatorname{diag}(s_1,s_2)$ with $s_1s_2\ne0$ by $\operatorname{diag}(1/s_1,1/s_2)$. For $2\times2$ matrices in general, \cref{prop:inverse-2x2} gives the inverse explicitly; for instance
\[
\begin{bmatrix}2&1\\1&2\end{bmatrix}^{-1}=\frac13\begin{bmatrix}2&-1\\-1&2\end{bmatrix}.
\]

\subsection{Singular Matrices}

A matrix fails to be invertible exactly when its map loses information, by \cref{prop:one-sided-inverses}\,(iv). The orthogonal projection $P_\alpha$ onto a line maps the whole orthogonal line to $\mathbf 0$, for example $P_\alpha(-\sin\alpha,\cos\alpha)=\mathbf 0$, and no map can recover a vector from its projection: the projection is singular. The same argument shows that a zero divisor is never invertible: if $AB=0$ with $B\ne0$, some column $\mathbf b$ of $B$ is nonzero and $A\mathbf b=\mathbf 0$, and an inverse would give $\mathbf b=A^{-1}A\mathbf b=\mathbf 0$. In particular the nilpotent matrix $N$ with $N^2=0$ is singular.

For a $2\times2$ matrix, singularity means $\det A=ad-bc=0$, by \cref{prop:inverse-2x2}, and for $n\times n$ matrices it means $\det A=0$, by \cref{thm:invertible-matrices}. Since $\det A$ is a polynomial in the entries, the singular matrices form a thin set, a hypersurface in the space $\mathbb R^{n\times n}$, and a matrix whose entries are drawn independently from a continuous probability distribution, such as the uniform distribution on $[0,1]$, is invertible with probability one. This does not make singular matrices irrelevant: matrices that are \emph{close} to singular are common in practice, and the inverse of such a matrix has huge entries. The matrix with the rows $(1,1)$ and $(1,1+\varepsilon)$, for example, has the determinant $\varepsilon$ and the inverse
\[
\frac1\varepsilon\begin{bmatrix}1+\varepsilon&-1\\-1&1\end{bmatrix},
\]
which blows up as $\varepsilon\to0$. How close a matrix is to being singular, relative to its size, is measured by its condition number, which \cref{sec:trace-norms} prepares and Chapter~21 and Volume~VIII develop.

If $A$ is invertible, the system $A\mathbf x=\mathbf b$ has the unique solution $\mathbf x=A^{-1}\mathbf b$. This formula is the right way to \emph{think} about the solution, but not the right way to \emph{compute} it: \cref{sec:computing-matrices} explains why the inverse is almost never formed in practice.

\section{Special Matrices}
\label{sec:special-matrices}

Matrices that arise in practice are rarely arbitrary tables of numbers. They have zeros in prescribed places, symmetries, or a structure inherited from the problem, and much of numerical linear algebra consists of recognising and exploiting such structure. This section collects the most important kinds. \Cref{tab:matrix-classes} at the end summarises which of them are preserved by products, inverses and transposition.

\subsection{Diagonal Matrices}

A \term{diagonal matrix} $D=\operatorname{diag}(d_1,\ldots,d_n)$ has the entries $d_i$ on the main diagonal and zeros elsewhere, $(D)_{ij}=d_i\delta_{ij}$. It is the simplest kind of matrix to compute with, because it acts on each coordinate separately.

\begin{proposition}[Diagonal matrices]
\label{prop:diagonal}
Let $D$ and $D'$ be diagonal matrices of order $m$ with the diagonal entries $d_i$ and $d_i'$, and let $A\in\mathbb R^{m\times n}$.
\begin{enumerate}
\item $DA$ is $A$ with its $i$th row multiplied by $d_i$, for every $i$; similarly, $A\operatorname{diag}(c_1,\ldots,c_n)$ is $A$ with its $j$th column multiplied by $c_j$.
\item $DD'=D'D=\operatorname{diag}(d_1d_1',\ldots,d_md_m')$.
\item $D$ is invertible if and only if every $d_i\ne0$, and then $D^{-1}=\operatorname{diag}(1/d_1,\ldots,1/d_m)$.
\end{enumerate}
\end{proposition}

\begin{proof}
(i) $(DA)_{ik}=\sum_jd_i\delta_{ij}a_{jk}=d_ia_{ik}$, and similarly $(A\operatorname{diag}(c_j))_{ik}=a_{ik}c_k$. (ii) is (i) with $A=D'$. (iii) If all $d_i\ne0$, the proposed inverse works by (ii). If $d_i=0$, then $D\mathbf e_i=\mathbf 0$, and $D$ is singular.
\end{proof}

Diagonal matrices describe changes of units. If the components of a vector are measured in different units, metres, seconds and kilograms, converting to other units multiplies each component by its own factor, that is, applies a diagonal matrix. Scaling the rows and columns of a matrix in this way, to make its entries of comparable size, is called \term{equilibration}, and it can improve the accuracy of computations (Volume~VIII).

\subsection{Triangular Matrices}

A square matrix $U$ is \term{upper triangular} if $u_{ij}=0$ for $i>j$, so that all entries below the main diagonal vanish, and \term{lower triangular} if all entries above the main diagonal vanish. It is \term{strictly} triangular if, in addition, its diagonal entries vanish, and \term{unit} triangular if its diagonal entries are all $1$. A matrix is lower triangular exactly when its transpose is upper triangular, so every statement about one kind has a mirror image for the other.

\begin{theorem}[Triangular matrices]
\label{thm:triangular}
Let $U,V\in\mathbb R^{n\times n}$ be upper triangular.
\begin{enumerate}
\item $UV$ is upper triangular, with the diagonal entries $(UV)_{ii}=u_{ii}v_{ii}$.
\item $U$ is invertible if and only if all its diagonal entries are nonzero, and then $U^{-1}$ is upper triangular with the diagonal entries $1/u_{ii}$.
\item Products and inverses of unit upper triangular matrices are unit upper triangular.
\end{enumerate}
The same statements hold for lower triangular matrices.
\end{theorem}

\begin{proof}
(i) In $(UV)_{ik}=\sum_ju_{ij}v_{jk}$ a term can be nonzero only if $i\le j$ and $j\le k$. For $i>k$ there is no such $j$, so the entry is zero; for $i=k$ the only possible $j$ is $i$, and the entry is $u_{ii}v_{ii}$.

(ii) Suppose first that all $u_{ii}\ne0$. For fixed $k$ we solve $U\mathbf x=\mathbf e_k$, whose $i$th equation is $u_{ii}x_i+\sum_{j>i}u_{ij}x_j=\delta_{ik}$, from the last equation upwards:
\[
x_i=\frac1{u_{ii}}\Bigl(\delta_{ik}-\sum_{j=i+1}^nu_{ij}x_j\Bigr)\qquad(i=n,n-1,\ldots,1).
\]
This \term{back substitution} determines $x_n,x_{n-1},\ldots,x_1$ in turn. For $i>k$ we get $x_i=0$, by downward induction, since $\delta_{ik}=0$ and the $x_j$ with $j>i$ are already zero; and then $x_k=1/u_{kk}$. So the matrix $X$ whose columns are these solutions is upper triangular, has the diagonal entries $1/u_{kk}$, and satisfies $UX=I$. The same construction applied to $X$, which is upper triangular with nonzero diagonal entries, gives an upper triangular $Y$ with $XY=I$. Then $U=U(XY)=(UX)Y=Y$, so $XU=I$ as well, and $X=U^{-1}$.

Conversely, suppose that some diagonal entry vanishes, and let $k$ be the smallest index with $u_{kk}=0$. We construct $\mathbf x\ne\mathbf 0$ with $U\mathbf x=\mathbf 0$, which shows that $U$ is singular. Put $x_k\coloneqq1$ and $x_j\coloneqq0$ for $j>k$. The upper left $(k-1)\times(k-1)$ part $U_{k-1}$ of $U$ is upper triangular with nonzero diagonal entries, hence invertible by the first part, and we choose $(x_1,\ldots,x_{k-1})\coloneqq-U_{k-1}^{-1}(u_{1k},\ldots,u_{k-1,k})$ (nothing to choose if $k=1$). Then the $i$th component of $U\mathbf x$ is, for $i<k$, $\sum_{j<k}u_{ij}x_j+u_{ik}=0$ by the choice of $x_1,\ldots,x_{k-1}$; for $i=k$, it is $u_{kk}x_k=0$, since $u_{kj}=0$ for $j<k$; and for $i>k$ it is $\sum_{j\ge i}u_{ij}x_j=0$, since $x_j=0$ for $j\ge i>k$.

(iii) follows from (i) and (ii), since $1\cdot1=1$ and $1/1=1$. For lower triangular matrices, apply (i) to (iii) to the transposes, using the reversal rules for transposes and inverses.
\end{proof}

For example,
\[
U=\begin{bmatrix}1&2&3\\0&1&4\\0&0&1\end{bmatrix},\qquad U^{-1}=\begin{bmatrix}1&-2&5\\0&1&-4\\0&0&1\end{bmatrix},
\]
as multiplying out confirms. Triangular matrices are central because triangular systems are easy: back substitution solves $U\mathbf x=\mathbf b$ with about $n^2$ operations (\cref{sec:computing-matrices}), and \cref{ch:linear-systems} reduces a general system to two triangular ones by writing $A=LU$ with $L$ unit lower triangular and $U$ upper triangular. Parts (i) and (iii) of the theorem are what make this factorisation work.

\subsection{Permutation Matrices}

A permutation of a set is a bijection of the set onto itself (\cref{sec:injections}), and the $n!$ permutations of $\{1,\ldots,n\}$ (\cref{sec:perm-comb}) rearrange the components of vectors. For a permutation $\sigma$ of $\{1,\ldots,n\}$, the \term{permutation matrix} $P_\sigma\in\mathbb R^{n\times n}$ is the matrix with the columns $\mathbf e_{\sigma(1)},\ldots,\mathbf e_{\sigma(n)}$; its entries are $1$ in the positions $(\sigma(j),j)$ and $0$ elsewhere, so it has exactly one $1$ in every row and every column. For $n=3$ and the permutation $\sigma$ with $\sigma(1)=2$, $\sigma(2)=3$, $\sigma(3)=1$,
\[
P_\sigma=\begin{bmatrix}0&0&1\\1&0&0\\0&1&0\end{bmatrix},\qquad P_\sigma\begin{bmatrix}x_1\\x_2\\x_3\end{bmatrix}=\begin{bmatrix}x_3\\x_1\\x_2\end{bmatrix}:
\]
the component $x_j$ moves to position $\sigma(j)$.

\begin{proposition}[Permutation matrices]
\label{prop:permutation-matrices}
Let $\sigma$ and $\tau$ be permutations of $\{1,\ldots,n\}$.
\begin{enumerate}
\item $P_\sigma\mathbf e_j=\mathbf e_{\sigma(j)}$, and $(P_\sigma\mathbf x)_i=x_{\sigma^{-1}(i)}$ for every $\mathbf x\in\mathbb R^n$. For a matrix $A$ with $n$ rows, $P_\sigma A$ is $A$ with its $j$th row moved to position $\sigma(j)$.
\item $P_\sigma P_\tau=P_{\sigma\circ\tau}$, and $P_{\operatorname{id}}=I$.
\item $P_\sigma$ is invertible, and $P_\sigma^{-1}=P_\sigma^{\mathsf T}=P_{\sigma^{-1}}$.
\end{enumerate}
\end{proposition}

\begin{proof}
(i) The first statement is the definition, by \cref{thm:matrices-linear-maps}. Then $P_\sigma\mathbf x=\sum_jx_j\mathbf e_{\sigma(j)}$, and the $i$th component of this sum is $x_j$ for the unique $j$ with $\sigma(j)=i$, that is, $j=\sigma^{-1}(i)$. Applied to each column of $A$, this moves the $j$th row to position $\sigma(j)$. (ii) $P_\sigma P_\tau\mathbf e_j=P_\sigma\mathbf e_{\tau(j)}=\mathbf e_{\sigma(\tau(j))}=P_{\sigma\circ\tau}\mathbf e_j$ for every $j$, so the two matrices have the same columns. (iii) By (ii), $P_\sigma P_{\sigma^{-1}}=P_{\operatorname{id}}=I=P_{\sigma^{-1}}P_\sigma$. The $(j,k)$ entry of $P_\sigma^{\mathsf T}P_\sigma$ is $\langle\mathbf e_{\sigma(j)},\mathbf e_{\sigma(k)}\rangle$, which is $1$ if $\sigma(j)=\sigma(k)$, that is, $j=k$, and $0$ otherwise, so $P_\sigma^{\mathsf T}P_\sigma=I$. By \cref{prop:one-sided-inverses}\,(i), the left inverse $P_\sigma^{\mathsf T}$ equals the inverse $P_{\sigma^{-1}}$.
\end{proof}

Multiplying by a permutation matrix from the right permutes the columns instead: $(AP_\sigma)_{ik}=a_{i\sigma(k)}$, so the $k$th column of $AP_\sigma$ is the $\sigma(k)$th column of $A$. The simplest permutation matrices exchange two rows; Gaussian elimination with pivoting (\cref{ch:linear-systems}) uses them to bring a suitable entry into the leading position, and writes $PA=LU$. Part~(ii) says that $\sigma\mapsto P_\sigma$ turns the composition of permutations into the product of matrices, so the permutation matrices form a group isomorphic to the group of permutations. In a computer, a permutation matrix is never stored as a matrix of zeros and ones: one stores the list $\sigma(1),\ldots,\sigma(n)$ and applies it by moving entries, which costs no arithmetic at all. Sorting an array, for instance, computes such a list (\cref{sec:injections}), and Volume~VI does this throughout.

\subsection{Orthogonal Matrices}

The rotations, the reflections and the permutation matrices share a property: they preserve lengths and angles.

\begin{theorem}[Orthogonal matrices]
\label{thm:orthogonal}
For $Q\in\mathbb R^{n\times n}$ the following are equivalent:
\begin{enumerate}
\item $Q^{\mathsf T}Q=I$;
\item the columns of $Q$ are orthonormal;
\item $\langle Q\mathbf x,Q\mathbf y\rangle=\langle\mathbf x,\mathbf y\rangle$ for all $\mathbf x,\mathbf y\in\mathbb R^n$;
\item $\|Q\mathbf x\|=\|\mathbf x\|$ for all $\mathbf x\in\mathbb R^n$.
\end{enumerate}
A matrix with these properties is \term{orthogonal}. It is invertible, $Q^{-1}=Q^{\mathsf T}$, and $Q^{\mathsf T}$ is orthogonal as well, so that the rows of $Q$ are also orthonormal. Products and inverses of orthogonal matrices are orthogonal.
\end{theorem}

\begin{proof}
(i)$\Leftrightarrow$(ii): the entries of $Q^{\mathsf T}Q$ are the inner products of the columns (\cref{prop:gram}). (i)$\Rightarrow$(iii): $\langle Q\mathbf x,Q\mathbf y\rangle=\langle\mathbf x,Q^{\mathsf T}Q\mathbf y\rangle=\langle\mathbf x,\mathbf y\rangle$ by \cref{prop:transpose-rules}\,(v). (iii)$\Rightarrow$(iv): take $\mathbf y=\mathbf x$. (iv)$\Rightarrow$(iii): by \eqref{eq:norm-expansion}, the inner product is determined by norms through the \term{polarisation identity}
\[
\langle\mathbf x,\mathbf y\rangle=\frac14\bigl(\|\mathbf x+\mathbf y\|^2-\|\mathbf x-\mathbf y\|^2\bigr),
\]
and $Q$ preserves the norms on the right, since $Q(\mathbf x\pm\mathbf y)=Q\mathbf x\pm Q\mathbf y$. (iii)$\Rightarrow$(i): the $(i,j)$ entry of $Q^{\mathsf T}Q$ is $\langle Q\mathbf e_i,Q\mathbf e_j\rangle=\langle\mathbf e_i,\mathbf e_j\rangle=\delta_{ij}$.

By (i) and \cref{thm:one-sided-square}, $Q^{\mathsf T}=Q^{-1}$, so $QQ^{\mathsf T}=I$ too, which is (i) for $Q^{\mathsf T}$. If $Q_1$ and $Q_2$ are orthogonal, then $(Q_1Q_2)^{\mathsf T}Q_1Q_2=Q_2^{\mathsf T}Q_1^{\mathsf T}Q_1Q_2=Q_2^{\mathsf T}Q_2=I$, and $Q_1^{-1}=Q_1^{\mathsf T}$ is orthogonal, as just shown.
\end{proof}

The rotations $R_\varphi$ and the reflections $S_\alpha$ are orthogonal, since their columns are orthonormal, and so are the permutation matrices, by \cref{prop:permutation-matrices}\,(iii). In the plane there are no others. The first column of an orthogonal $2\times2$ matrix is a unit vector, $(\cos\varphi,\sin\varphi)$ for some $\varphi$ (\cref{prop:circle-points}), and the second is a unit vector orthogonal to it, of which there are exactly two, $\pm(-\sin\varphi,\cos\varphi)$; the plus sign gives $R_\varphi$ and the minus sign the reflection $S_{\varphi/2}$. In higher dimensions every orthogonal matrix is a product of at most $n$ reflections of the kind described next, as \cref{ch:inner-products} shows.

Orthogonal matrices are the matrices of rigid motions that fix the origin, and they are the safest matrices to compute with. Multiplying by $Q$ does not change lengths, so it cannot magnify errors already present in a vector; and the inverse is available for free, as the transpose. For this reason the most reliable algorithms of numerical linear algebra transform a problem by orthogonal matrices only (\cref{ch:inner-products}, Volume~VIII). The workhorse among them is the following reflection.

\begin{proposition}[Householder reflections]
\label{prop:householder}
Let $\mathbf u\in\mathbb R^n$ be a unit vector and $H\coloneqq I-2\mathbf u\mathbf u^{\mathsf T}$. Then:
\begin{enumerate}
\item $H$ is symmetric and orthogonal, and $H^2=I$, so that $H^{-1}=H$;
\item $H\mathbf u=-\mathbf u$, and $H\mathbf x=\mathbf x$ for every $\mathbf x\perp\mathbf u$;
\item if $\mathbf x\ne\mathbf y$ are vectors with $\|\mathbf x\|=\|\mathbf y\|$ and $\mathbf u\coloneqq(\mathbf x-\mathbf y)/\|\mathbf x-\mathbf y\|$, then $H\mathbf x=\mathbf y$.
\end{enumerate}
\end{proposition}

\begin{proof}
(i) $H^{\mathsf T}=I-2(\mathbf u\mathbf u^{\mathsf T})^{\mathsf T}=H$ by \cref{prop:transpose-rules}\,(iv), and $H^2=I-4\mathbf u\mathbf u^{\mathsf T}+4\mathbf u(\mathbf u^{\mathsf T}\mathbf u)\mathbf u^{\mathsf T}=I$, since $\mathbf u^{\mathsf T}\mathbf u=1$; so $H^{\mathsf T}H=H^2=I$. (ii) $H\mathbf x=\mathbf x-2\mathbf u\langle\mathbf u,\mathbf x\rangle$, which is $\mathbf u-2\mathbf u=-\mathbf u$ for $\mathbf x=\mathbf u$ and $\mathbf x$ for $\mathbf x\perp\mathbf u$. (iii) Since $\langle\mathbf x-\mathbf y,\mathbf x+\mathbf y\rangle=\|\mathbf x\|^2-\|\mathbf y\|^2=0$, the vector $\mathbf x+\mathbf y$ is orthogonal to $\mathbf u$. Writing $\mathbf x=(\mathbf x-\mathbf y)/2+(\mathbf x+\mathbf y)/2$, (ii) gives $H\mathbf x=-(\mathbf x-\mathbf y)/2+(\mathbf x+\mathbf y)/2=\mathbf y$.
\end{proof}

The matrix $H$ is the reflection in the hyperplane $\{\mathbf x: \langle\mathbf u,\mathbf x\rangle=0\}$: it keeps the hyperplane fixed and reverses the normal direction, as the reflections $S_\alpha$ do in the plane. Part~(iii) says that any vector can be reflected onto any other vector of the same length, and in particular onto a multiple of $\mathbf e_1$. For $\mathbf x=(3,4)$ and $\mathbf y=(5,0)$, the vector $\mathbf x-\mathbf y=(-2,4)$ has the squared length $20$, and $\langle\mathbf x-\mathbf y,\mathbf x\rangle=10$, and so $H\mathbf x=\mathbf x-2(\mathbf x-\mathbf y)\langle\mathbf x-\mathbf y,\mathbf x\rangle/20=(3,4)-(-2,4)=(5,0)$. Reflecting the columns of a matrix in this way, one after the other, introduces zeros below the diagonal while preserving lengths; this is Householder's method for the QR factorisation~\cite{householder1958}, which \cref{ch:inner-products} and Volume~VIII develop, and where one chooses the sign of $\mathbf y=\pm\|\mathbf x\|\mathbf e_1$ to avoid cancellation in $\mathbf x-\mathbf y$. In practice $H$ is never formed: $H\mathbf x=\mathbf x-2\mathbf u\langle\mathbf u,\mathbf x\rangle$ costs about $4n$ operations, whereas the matrix would take $n^2$ numbers to store and $2n^2$ operations to apply.

\subsection{Symmetric, Idempotent and Nilpotent Matrices}

Symmetric matrices were introduced in \cref{sec:transpose}, and the positive definite ones, symmetric matrices with $\langle\mathbf x,A\mathbf x\rangle>0$ for all $\mathbf x\ne\mathbf 0$, in \cref{def:definite}. A positive definite matrix is invertible, by \cref{thm:invertible-matrices}\,(i), since $A\mathbf x=\mathbf 0$ implies $\langle\mathbf x,A\mathbf x\rangle=0$ and hence $\mathbf x=\mathbf 0$. \Cref{ch:symmetric-matrices} develops their theory; for now we note that the Gram matrices $A^{\mathsf T}A$ of \cref{prop:gram} are positive semidefinite, and positive definite when $A$ maps only $\mathbf 0$ to $\mathbf 0$.

A square matrix $P$ is \term{idempotent} if $P^2=P$: applying it twice is the same as applying it once. The orthogonal projection $P_\alpha=\mathbf u\mathbf u^{\mathsf T}$ onto a line, with a unit vector $\mathbf u$, is idempotent, since $(\mathbf u\mathbf u^{\mathsf T})^2=\mathbf u(\mathbf u^{\mathsf T}\mathbf u)\mathbf u^{\mathsf T}=\mathbf u\mathbf u^{\mathsf T}$, and so is the \term{oblique projection}
\[
P=\begin{bmatrix}1&1\\0&0\end{bmatrix},\qquad P\begin{bmatrix}x_1\\x_2\end{bmatrix}=\begin{bmatrix}x_1+x_2\\0\end{bmatrix},
\]
which maps the plane onto the $x_1$-axis along the lines $x_1+x_2=c$, in the direction $(1,-1)$ rather than orthogonally. Every idempotent $P$ splits vectors into two parts, $\mathbf x=P\mathbf x+(I-P)\mathbf x$: the first part is fixed by $P$, since $P(P\mathbf x)=P\mathbf x$, and the second is mapped to $\mathbf 0$, since $P(I-P)=P-P^2=0$. If $P$ is moreover symmetric, the two parts are orthogonal, $\langle P\mathbf x,(I-P)\mathbf y\rangle=\langle\mathbf x,(P-P^2)\mathbf y\rangle=0$, and $P$ is an \term{orthogonal projection}; \cref{ch:inner-products} shows that every symmetric idempotent matrix is the orthogonal projection onto a subspace. The identity is the only invertible idempotent matrix, since $P^2=P$ gives $P=P^{-1}P^2=P^{-1}P=I$.

A square matrix $N$ is \term{nilpotent} if $N^k=0$ for some $k\ge1$. The matrix $N$ with $N^2=0$ of \cref{sec:matrix-multiplication} is an example, and so is the differentiation matrix $D$ of \cref{sec:matrices-linear-maps}: differentiating a polynomial of degree at most $3$ four times gives zero, so $D^4=0$. Both are strictly upper triangular, and this is the general pattern.

\begin{proposition}[Nilpotent matrices]
\label{prop:nilpotent}
Let $N\in\mathbb R^{n\times n}$ be strictly upper triangular. Then $(N^j)_{ik}=0$ whenever $k-i<j$; in particular $N^n=0$. If $N^k=0$, then $I-N$ is invertible, and
\[
(I-N)^{-1}=I+N+N^2+\cdots+N^{k-1}.
\]
\end{proposition}

\begin{proof}
By induction on $j$: for $j=1$ the statement is the definition of a strictly upper triangular matrix. If it holds for $j$, a term $(N^j)_{il}n_{lk}$ of $(N^{j+1})_{ik}$ can be nonzero only if $l-i\ge j$ and $k-l\ge1$, so only if $k-i\ge j+1$. Since $k-i\le n-1$ for all entries, $N^n=0$. Finally, the finite geometric sum gives $(I-N)(I+N+\cdots+N^{k-1})=I-N^k=I$, and the factors commute, so the product in the other order is $I$ as well.
\end{proof}

The proposition gives another way to invert a unit upper triangular matrix. The matrix $U$ of the example after \cref{thm:triangular} is $I-N$ with
\[
N=\begin{bmatrix}0&-2&-3\\0&0&-4\\0&0&0\end{bmatrix},\qquad N^2=\begin{bmatrix}0&0&8\\0&0&0\\0&0&0\end{bmatrix},\qquad N^3=0,
\]
and $U^{-1}=I+N+N^2$ is the matrix computed there. The formula is the finite version of the geometric series $1/(1-x)=\sum_kx^k$, and \cref{sec:trace-norms} extends it to infinite series of matrices.

\subsection{Banded and Sparse Matrices}

Many matrices that arise from differential equations, from splines or from time series have their nonzero entries concentrated near the diagonal. A matrix $A$ has \term{lower bandwidth} $p$ and \term{upper bandwidth} $q$ if $a_{ij}=0$ whenever $i-j>p$ or $j-i>q$. Diagonal matrices have the bandwidths $(0,0)$, upper triangular matrices $(0,n-1)$, \term{tridiagonal} matrices $(1,1)$ and \term{bidiagonal} matrices $(0,1)$ or $(1,0)$; the difference matrix $\Delta$ of \cref{sec:matrices-linear-maps} is an upper bidiagonal matrix that is not square. A banded matrix of order $n$ has at most $(p+q+1)n$ nonzero entries, far fewer than $n^2$ when the bandwidths are small, and it is stored and multiplied in proportion to this number.

\begin{proposition}[Products of banded matrices]
\label{prop:band-product}
If $A$ has the bandwidths $(p_1,q_1)$ and $B$ has the bandwidths $(p_2,q_2)$, then $AB$ has lower bandwidth at most $p_1+p_2$ and upper bandwidth at most $q_1+q_2$.
\end{proposition}

\begin{proof}
A term $a_{ij}b_{jk}$ of $(AB)_{ik}$ can be nonzero only if $i-j\le p_1$ and $j-k\le p_2$, which implies $i-k\le p_1+p_2$, and only if $j-i\le q_1$ and $k-j\le q_2$, which implies $k-i\le q_1+q_2$.
\end{proof}

The most important tridiagonal matrix is the \term{second difference matrix}
\[
T_n=\begin{bmatrix}2&-1&&&\\-1&2&-1&&\\&\ddots&\ddots&\ddots&\\&&-1&2&-1\\&&&-1&2\end{bmatrix}\in\mathbb R^{n\times n},
\]
with $2$ on the diagonal, $-1$ next to it and zeros elsewhere (blank entries in a displayed matrix are zero). It comes from the second difference quotient of \cref{prop:difference-taylor}\,(iii). Divide $[0,1]$ into $n+1$ intervals of length $h=1/(n+1)$, with the grid points $x_i=ih$, and let $u$ be a smooth function with $u(0)=u(1)=0$. With $u_i\coloneqq u(x_i)$, the $i$th component of $T_n\mathbf u/h^2$ is $-(u_{i-1}-2u_i+u_{i+1})/h^2=-u''(x_i)+O(h^2)$, where $u_0=u_{n+1}=0$ are the boundary values. So $T_n/h^2$ is the matrix that approximates the operator $u\mapsto-u''$, and the boundary value problem $-u''=f$, $u(0)=u(1)=0$, becomes the linear system $T_n\mathbf u=h^2\mathbf f$, which Part~IV and Volume~XI solve.

The matrix $T_n$ has a remarkable amount of structure. It is the Gram matrix $\Delta_0^{\mathsf T}\Delta_0$ of the $(n+1)\times n$ difference matrix $\Delta_0$ that maps $\mathbf u$ to the differences $(u_1-u_0,u_2-u_1,\ldots,u_{n+1}-u_n)$ with $u_0=u_{n+1}=0$, as one checks by computing the inner products of its columns, so it is symmetric and positive definite (\cref{prop:gram}):
\[
\langle\mathbf u,T_n\mathbf u\rangle=\|\Delta_0\mathbf u\|^2=u_1^2+\sum_{i=1}^{n-1}(u_{i+1}-u_i)^2+u_n^2 .
\]
The vectors $\mathbf s$ with the components $s_i=\sin(i\theta)$, $\theta=\pi/(n+1)$, are mapped to multiples of themselves: by the addition formula $\sin((i-1)\theta)+\sin((i+1)\theta)=2\sin(i\theta)\cos\theta$, and since $s_0=s_{n+1}=0$, we get $T_n\mathbf s=2(1-\cos\theta)\mathbf s$, which makes $\mathbf s$ an eigenvector in the sense of \cref{ch:eigenvalues}. And its inverse can be written down.

\begin{proposition}[The inverse of the second difference matrix]
\label{prop:second-difference-inverse}
$T_n$ is invertible, and
\[
(T_n^{-1})_{ij}=\frac{\min(i,j)\,\bigl(n+1-\max(i,j)\bigr)}{n+1}\qquad(1\le i,j\le n).
\]
\end{proposition}

\begin{proof}
Let $G$ be the matrix with these entries. It suffices to show $T_nG=I$: since $T_n$ and $G$ are symmetric, the reversal rule then gives $GT_n=G^{\mathsf T}T_n^{\mathsf T}=(T_nG)^{\mathsf T}=I$ as well. Fix a column $j$ and write $g_i\coloneqq G_{ij}$, also for $i=0$ and $i=n+1$, where the formula gives $g_0=g_{n+1}=0$. The $i$th component of $T_n$ times the $j$th column of $G$ is $-g_{i-1}+2g_i-g_{i+1}$. For $0\le i\le j$, $g_i=i(n+1-j)/(n+1)$ is a linear function of $i$, and for $j\le i\le n+1$, $g_i=j(n+1-i)/(n+1)$ is another; the second difference of a linear function is zero, so the component vanishes for $i\ne j$. For $i=j$ it is
\[
\frac{2j(n+1-j)-(j-1)(n+1-j)-j(n-j)}{n+1}=\frac{(n+1-j)(j+1)-j(n-j)}{n+1}=1 ,
\]
since $(n+1-j)(j+1)-j(n-j)=n+1$. So the $j$th column of $T_nG$ is $\mathbf e_j$.
\end{proof}

For $n=3$,
\[
T_3=\begin{bmatrix}2&-1&0\\-1&2&-1\\0&-1&2\end{bmatrix},\qquad T_3^{-1}=\frac14\begin{bmatrix}3&2&1\\2&4&2\\1&2&3\end{bmatrix}.
\]
The inverse of the tridiagonal matrix $T_n$ has no zero entries at all. This is typical: the inverse of a banded or otherwise sparse matrix is usually full, and forming it would destroy the very structure that makes the matrix cheap to store and to use. It is one more reason to solve systems without forming inverses; \cref{ch:linear-systems} solves systems with $T_n$ in a number of operations proportional to $n$.

Matrices with only a small fraction of nonzero entries, whether banded or not, are called \term{sparse}. The adjacency matrix of a road network or of the links between web pages, the matrices of finite element methods and the data matrices of text documents are sparse, with perhaps millions of rows but only a handful of nonzero entries in each. Such matrices are stored as lists of their nonzero entries with their positions, and their algorithms are designed to create as few new nonzero entries as possible; Volume~IX is devoted to them.

\begin{table}
\sdfloat{%
\begin{sdtable}{colspec={llll}, row{1}={font=\small\bfseries}}
\toprule
kind & product & inverse (if it exists) & transpose \\
\midrule
diagonal & diagonal & diagonal & diagonal \\
upper triangular & upper triangular & upper triangular & lower triangular \\
unit upper triangular & unit upper triangular & unit upper triangular & unit lower triangular \\
permutation & permutation & permutation & permutation \\
orthogonal & orthogonal & orthogonal & orthogonal \\
symmetric & only if $AB=BA$ & symmetric & symmetric \\
tridiagonal & bandwidths $(2,2)$ & full in general & tridiagonal \\
idempotent & not in general & only $I$ & idempotent \\
nilpotent & not in general & never invertible & nilpotent \\
\bottomrule
\end{sdtable}}{Kinds of square matrices of the same order and what becomes of them under products, inversion and transposition. ``Not in general'' means that the product of two matrices of the kind can fail to be of the kind. The entries follow from \cref{prop:diagonal,thm:triangular,prop:permutation-matrices,thm:orthogonal,prop:band-product,prop:second-difference-inverse} and the discussion of this section.\label{tab:matrix-classes}}
\end{table}

\Cref{tab:matrix-classes} summarises the behaviour of these kinds under the basic operations. A few entries deserve a comment. The product of two symmetric matrices is symmetric exactly when they commute, since $(AB)^{\mathsf T}=B^{\mathsf T}A^{\mathsf T}=BA$; for example $\operatorname{diag}(1,2)$ times the matrix with the rows $(0,1)$ and $(1,0)$ is not symmetric. The inverse of a symmetric matrix is symmetric, by \cref{prop:inverse-rules}\,(iii). The projections $\operatorname{diag}(1,0)$ and $\mathbf u\mathbf u^{\mathsf T}$ with $\mathbf u=(1,1)/\sqrt2$ are idempotent, but their product, with the rows $(1/2,1/2)$ and $(0,0)$, squares to half of itself. And the nilpotent matrices $N$ and $N^{\mathsf T}$ of \cref{sec:matrix-multiplication} have the product $\operatorname{diag}(1,0)$, which is idempotent and not nilpotent.

\section{Block Matrices}
\label{sec:block-matrices}

\subsection{Partitioned Matrices and Block Multiplication}

A matrix can be cut by horizontal and vertical lines into smaller matrices, its \term{blocks}. For example, a $5\times5$ matrix cut after the second row and the second column is written
\[
M=\begin{bmatrix}A&B\\C&D\end{bmatrix},\qquad A\in\mathbb R^{2\times2},\ B\in\mathbb R^{2\times3},\ C\in\mathbb R^{3\times2},\ D\in\mathbb R^{3\times3}.
\]
We have already met partitioned matrices without the name: a matrix written as a row of its columns, $A=(\mathbf a_1\ \cdots\ \mathbf a_n)$, is partitioned into blocks of one column each, and the affine map in homogeneous coordinates in \cref{sec:matrices-linear-maps} was a $2\times2$ arrangement of blocks of different sizes. The point of partitioning is that blocks can be multiplied as if they were numbers, provided that the cuts match and the order of the factors is kept.

\begin{theorem}[Block multiplication]
\label{thm:block-multiplication}
Let $A\in\mathbb R^{m\times n}$ and $B\in\mathbb R^{n\times p}$ be partitioned as
\[
A=\begin{bmatrix}A_{11}&\cdots&A_{1s}\\\vdots&&\vdots\\A_{r1}&\cdots&A_{rs}\end{bmatrix},\qquad B=\begin{bmatrix}B_{11}&\cdots&B_{1t}\\\vdots&&\vdots\\B_{s1}&\cdots&B_{st}\end{bmatrix},
\]
where $A_{IJ}\in\mathbb R^{m_I\times n_J}$ and $B_{JK}\in\mathbb R^{n_J\times p_K}$; that is, the columns of $A$ and the rows of $B$ are cut in the same places. Then $AB$, partitioned into rows of heights $m_1,\ldots,m_r$ and columns of widths $p_1,\ldots,p_t$, has the blocks
\[
(AB)_{IK}=\sum_{J=1}^sA_{IJ}B_{JK}\qquad(1\le I\le r,\ 1\le K\le t).
\]
\end{theorem}

\begin{proof}
Let the entry $(i,k)$ of $AB$ lie in the block $(I,K)$. The sum $(AB)_{ik}=\sum_{j=1}^na_{ij}b_{jk}$ can be split into $s$ partial sums, one over the indices $j$ of each column group $J$ of $A$, which is also a row group of $B$. The partial sum for $J$ runs over the entries of row $i$ of $A$ that lie in $A_{IJ}$ and the entries of column $k$ of $B$ that lie in $B_{JK}$, so it is the corresponding entry of the product $A_{IJ}B_{JK}$. Adding over $J$ gives the claim.
\end{proof}

For two blocks in each direction,
\[
\begin{bmatrix}A&B\\C&D\end{bmatrix}\begin{bmatrix}E&F\\G&H\end{bmatrix}=\begin{bmatrix}AE+BG&AF+BH\\CE+DG&CF+DH\end{bmatrix},
\]
the formula for $2\times2$ matrices with the order of every product preserved. The four readings of \cref{thm:four-readings} are special cases: cutting $B$ into columns gives the column reading, cutting $A$ into rows the row reading, and cutting $A$ into columns and $B$ into rows the outer product reading. A \term{block diagonal} matrix $\operatorname{diag}(A_1,\ldots,A_k)$, with square blocks $A_i$ on the diagonal and zero blocks elsewhere, acts on each group of coordinates separately, and by the theorem the product of two block diagonal matrices with the same block sizes is block diagonal with the products $A_iB_i$ as blocks.

Block multiplication is more than a notational convenience. It allows an algorithm to work on blocks that fit into the fast memory of a computer, which is how high-performance libraries multiply matrices (\cref{sec:matrix-multiplication}); it allows a large computation to be distributed among many processors, each of which owns some blocks; and it is the basis of recursive algorithms, which split a matrix into four blocks and treat each block in the same way, like Strassen's algorithm in the note at the end of this section.

\subsection{Block Triangular Matrices and the Schur Complement}

Elimination, the subject of \cref{ch:linear-systems}, removes unknowns from a system one at a time. Done with blocks, it removes a whole group of unknowns at once, and the matrix that remains has a name.

\begin{theorem}[Block elimination and the Schur complement]
\label{thm:schur-complement}
Let
\[
M=\begin{bmatrix}A&B\\C&D\end{bmatrix}
\]
with $A\in\mathbb R^{k\times k}$ invertible and $D\in\mathbb R^{l\times l}$, and let $S\coloneqq D-CA^{-1}B\in\mathbb R^{l\times l}$ be the \term{Schur complement} of $A$ in $M$.
\begin{enumerate}
\item $M$ has the block $LDU$ factorisation
\[
M=\begin{bmatrix}I&0\\CA^{-1}&I\end{bmatrix}\begin{bmatrix}A&0\\0&S\end{bmatrix}\begin{bmatrix}I&A^{-1}B\\0&I\end{bmatrix}.
\]
\item $M$ is invertible if and only if $S$ is invertible, and then
\[
M^{-1}=\begin{bmatrix}A^{-1}+A^{-1}BS^{-1}CA^{-1}&-A^{-1}BS^{-1}\\-S^{-1}CA^{-1}&S^{-1}\end{bmatrix}.
\]
\end{enumerate}
\end{theorem}

\begin{proof}
(i) By block multiplication, the first two factors give the matrix with the block rows $(A,0)$ and $(C,S)$, and multiplying this by the third factor gives the block rows $(A,B)$ and $(C,CA^{-1}B+S)=(C,D)$.

(ii) The outer factors of (i) are invertible, and their inverses are obtained by changing the sign of the off-diagonal block: for instance
\[
\begin{bmatrix}I&0\\CA^{-1}&I\end{bmatrix}\begin{bmatrix}I&0\\-CA^{-1}&I\end{bmatrix}=\begin{bmatrix}I&0\\CA^{-1}-CA^{-1}&I\end{bmatrix}=I,
\]
and similarly in the other order and for the third factor. If $S$ is invertible, the middle factor $\Delta\coloneqq\operatorname{diag}(A,S)$ has the inverse $\operatorname{diag}(A^{-1},S^{-1})$, by block multiplication, so $M$ is a product of invertible matrices and is invertible, with
\[
M^{-1}=\begin{bmatrix}I&-A^{-1}B\\0&I\end{bmatrix}\begin{bmatrix}A^{-1}&0\\0&S^{-1}\end{bmatrix}\begin{bmatrix}I&0\\-CA^{-1}&I\end{bmatrix}
\]
by \cref{prop:inverse-rules}\,(ii); multiplying out gives the formula. Conversely, if $M$ is invertible, so is $\Delta$, as a product of $M$ and the inverses of the outer factors. Partition $X\coloneqq\Delta^{-1}$ like $\Delta$. The lower right blocks of $\Delta X=I$ and $X\Delta=I$ give $SX_{22}=I$ and $X_{22}S=I$, so $S$ is invertible.
\end{proof}

The factorisation (i) records block elimination. Multiplying $M$ from the left by the inverse of the first factor subtracts $CA^{-1}$ times the first block row from the second, which clears the block $C$ and leaves $S$ in the corner:
\[
\begin{bmatrix}I&0\\-CA^{-1}&I\end{bmatrix}\begin{bmatrix}A&B\\C&D\end{bmatrix}=\begin{bmatrix}A&B\\0&S\end{bmatrix}.
\]
In terms of the system $A\mathbf x+B\mathbf y=\mathbf f$, $C\mathbf x+D\mathbf y=\mathbf g$, solving the first equation for $\mathbf x=A^{-1}(\mathbf f-B\mathbf y)$ and substituting into the second leaves the smaller system $S\mathbf y=\mathbf g-CA^{-1}\mathbf f$ for $\mathbf y$ alone. For $k=l=1$ the Schur complement of the number $a\ne0$ is $d-cb/a$, the entry that one step of Gaussian elimination leaves behind, and \cref{ch:linear-systems} obtains the LU factorisation by repeating this step. The determinant of $M$ is $\det A\det S$, by the factorisation and the product rule for determinants (\cref{ch:determinants}), and when $M$ is symmetric and positive definite, so is $S$ (\cref{ch:symmetric-matrices}); both facts are used often. The name was given by Haynsworth in 1968 in reference to a determinant lemma of Schur~\cite{schur1917}; Zhang's collection~\cite{zhang2005schur} surveys the history and the many uses of the notion.

With $C=0$ the theorem gives the inverse of a \term{block upper triangular} matrix: if $A$ and $D$ are invertible, then
\[
\begin{bmatrix}A&B\\0&D\end{bmatrix}^{-1}=\begin{bmatrix}A^{-1}&-A^{-1}BD^{-1}\\0&D^{-1}\end{bmatrix},
\]
since $S=D$; with $B=C=0$ as well, the inverse of a block diagonal matrix is block diagonal with the inverses of the blocks. These are the block versions of \cref{thm:triangular,prop:diagonal}.

A typical use is a \term{bordered matrix}, a matrix enlarged by one row and one column. Suppose that the inverse of $A\in\mathbb R^{k\times k}$ is known and that
\[
M=\begin{bmatrix}A&\mathbf b\\\mathbf c^{\mathsf T}&d\end{bmatrix}
\]
with $\mathbf b,\mathbf c\in\mathbb R^k$ and a number $d$. The Schur complement is the number $s=d-\mathbf c^{\mathsf T}A^{-1}\mathbf b$, and $M$ is invertible exactly when $s\ne0$; the theorem then gives $M^{-1}$ from $A^{-1}$ with about $4k^2$ operations instead of the $k^3$ or so of a fresh inversion. For
\[
A=\begin{bmatrix}2&1\\1&2\end{bmatrix},\quad\mathbf b=\begin{bmatrix}1\\0\end{bmatrix},\quad\mathbf c=\begin{bmatrix}0\\1\end{bmatrix},\quad d=3,\qquad M^{-1}=\begin{bmatrix}3/5&-1/5&-1/5\\-3/10&3/5&1/10\\1/10&-1/5&3/10\end{bmatrix},
\]
we have $A^{-1}\mathbf b=(2,-1)/3$, $\mathbf c^{\mathsf T}A^{-1}=(-1,2)/3$ and $s=3-(-1/3)=10/3$, so the lower right entry of $M^{-1}$ is $1/s=3/10$, the last column above it is $-A^{-1}\mathbf b/s=(-1/5,1/10)$, and the upper left block is $A^{-1}+(A^{-1}\mathbf b)(\mathbf c^{\mathsf T}A^{-1})/s$. Adding one row and column at a time in this way builds up inverses and factorisations recursively, and it is how many updating algorithms of statistics and signal processing work.

\subsection{Rank-One Updates: the Sherman--Morrison--Woodbury Formula}

Suppose again that $A^{-1}$ is known, and that $A$ changes slightly, by the addition of a rank-one matrix $\mathbf u\mathbf v^{\mathsf T}$: a single entry changes ($\mathbf u$ and $\mathbf v$ multiples of standard basis vectors), or a single row or column. The inverse of the new matrix need not be computed from scratch.

\begin{theorem}[Sherman--Morrison formula]
\label{thm:sherman-morrison}
Let $A\in\mathbb R^{n\times n}$ be invertible and $\mathbf u,\mathbf v\in\mathbb R^n$. Then $A+\mathbf u\mathbf v^{\mathsf T}$ is invertible if and only if $1+\mathbf v^{\mathsf T}A^{-1}\mathbf u\ne0$, and then
\[
\bigl(A+\mathbf u\mathbf v^{\mathsf T}\bigr)^{-1}=A^{-1}-\frac{A^{-1}\mathbf u\,\mathbf v^{\mathsf T}A^{-1}}{1+\mathbf v^{\mathsf T}A^{-1}\mathbf u}.
\]
\end{theorem}

\begin{proof}
Write $\beta\coloneqq1+\mathbf v^{\mathsf T}A^{-1}\mathbf u$, a number, and suppose $\beta\ne0$. Multiplying out, and noting that $\mathbf v^{\mathsf T}A^{-1}\mathbf u=\beta-1$ is a number that can be moved to the front,
\[
\bigl(A+\mathbf u\mathbf v^{\mathsf T}\bigr)\Bigl(A^{-1}-\frac{A^{-1}\mathbf u\mathbf v^{\mathsf T}A^{-1}}\beta\Bigr)=I+\mathbf u\mathbf v^{\mathsf T}A^{-1}-\frac{\mathbf u\mathbf v^{\mathsf T}A^{-1}+(\beta-1)\mathbf u\mathbf v^{\mathsf T}A^{-1}}\beta=I,
\]
and by \cref{thm:one-sided-square} the matrix in the second factor is the inverse. Conversely, if $\beta=0$, then $\mathbf u\ne\mathbf 0$ (otherwise $\beta=1$), so $\mathbf x\coloneqq A^{-1}\mathbf u\ne\mathbf 0$, and $(A+\mathbf u\mathbf v^{\mathsf T})\mathbf x=\mathbf u+\mathbf u(\mathbf v^{\mathsf T}A^{-1}\mathbf u)=\beta\mathbf u=\mathbf 0$, so $A+\mathbf u\mathbf v^{\mathsf T}$ is singular.
\end{proof}

For $A=\operatorname{diag}(1,2,4)$ and $\mathbf u=\mathbf v=(1,1,1)$, the formula needs only $A^{-1}\mathbf u=(1,1/2,1/4)$ and $\beta=1+1+1/2+1/4=11/4$:
\[
\bigl(A+\mathbf u\mathbf u^{\mathsf T}\bigr)^{-1}=\begin{bmatrix}1&0&0\\0&1/2&0\\0&0&1/4\end{bmatrix}-\frac4{11}\begin{bmatrix}1\\1/2\\1/4\end{bmatrix}\begin{bmatrix}1&1/2&1/4\end{bmatrix}=\frac1{22}\begin{bmatrix}14&-4&-2\\-4&9&-1\\-2&-1&5\end{bmatrix}.
\]
The cost of the update is a few matrix--vector products and one outer product, about $3n^2$ multiplications, instead of the $n^3$ or so of a new inversion, and if only the solution of a system is needed, even the outer product can be avoided: $(A+\mathbf u\mathbf v^{\mathsf T})^{-1}\mathbf b=\mathbf z-\mathbf w\,(\mathbf v^{\mathsf T}\mathbf z)/\beta$ with $\mathbf z=A^{-1}\mathbf b$ and $\mathbf w=A^{-1}\mathbf u$. Such updates occur in the quasi-Newton methods of optimisation (Volume~XII), where an approximate Hessian is corrected by rank-one or rank-two matrices in every step, and in statistics, where deleting one observation from a regression changes $X^{\mathsf T}X$ by a rank-one matrix (Volume~XV). Their numerical reliability is a separate question, which Volume~VIII discusses; Hager's survey~\cite{hager1989} describes the formula and its applications.

The formula generalises from rank one to rank $k$. A heuristic derivation shows where it comes from. For the matrix
\[
M=\begin{bmatrix}A&U\\V&-C^{-1}\end{bmatrix}
\]
with invertible $A\in\mathbb R^{n\times n}$ and $C\in\mathbb R^{k\times k}$, $U\in\mathbb R^{n\times k}$ and $V\in\mathbb R^{k\times n}$, \cref{thm:schur-complement} with the pivot block $A$ gives the upper left block of $M^{-1}$ as $A^{-1}-A^{-1}U(C^{-1}+VA^{-1}U)^{-1}VA^{-1}$. Eliminating with the other diagonal block $-C^{-1}$ as pivot instead, the same theorem with the roles of the blocks exchanged gives this block as the inverse of the Schur complement $A-U(-C^{-1})^{-1}V=A+UCV$. Equating the two expressions gives the following identity, which we now verify directly.

\begin{theorem}[Woodbury identity]
\label{thm:woodbury}
Let $A\in\mathbb R^{n\times n}$ and $C\in\mathbb R^{k\times k}$ be invertible, $U\in\mathbb R^{n\times k}$, $V\in\mathbb R^{k\times n}$, and suppose that $C^{-1}+VA^{-1}U$ is invertible. Then $A+UCV$ is invertible, and
\[
(A+UCV)^{-1}=A^{-1}-A^{-1}U\bigl(C^{-1}+VA^{-1}U\bigr)^{-1}VA^{-1}.
\]
\end{theorem}

\begin{proof}
Write $W\coloneqq C^{-1}+VA^{-1}U$. Then
\begin{align*}
(A+UCV)\bigl(A^{-1}-A^{-1}UW^{-1}VA^{-1}\bigr)&=I+UCVA^{-1}-\bigl(U+UCVA^{-1}U\bigr)W^{-1}VA^{-1}\\
&=I+UCVA^{-1}-UC\bigl(C^{-1}+VA^{-1}U\bigr)W^{-1}VA^{-1}\\
&=I+UCVA^{-1}-UCVA^{-1}=I,
\end{align*}
where the first line multiplies out and collects the terms ending in $W^{-1}VA^{-1}$, the second writes $U=UCC^{-1}$, and the third uses $WW^{-1}=I$. By \cref{thm:one-sided-square} the second factor is the inverse of $A+UCV$.
\end{proof}

With $k=1$, $C=1$, $U=\mathbf u$ and $V=\mathbf v^{\mathsf T}$ the identity is the Sherman--Morrison formula. Its value lies in the sizes: when $k$ is much smaller than $n$, the inverse of the $n\times n$ matrix $A+UCV$ is obtained from that of $A$ and the inverse of the small $k\times k$ matrix $W$. If $A$ is diagonal, for instance, as for a covariance matrix that is a diagonal matrix plus a low-rank part, a system with $A+UCV$ can be solved with a number of operations proportional to $nk^2$ instead of $n^3$. The formulas were found several times; Sherman and Morrison published the rank-one case for a change of a single entry~\cite{sherman1950}, Woodbury the general one in a report of the same year~\cite{woodbury1950}, and Henderson and Searle trace the history further back~\cite{henderson1981}.

\begin{note}{Strassen's algorithm and the exponent of matrix multiplication}
Block multiplication of two matrices cut into $2\times2$ blocks of order $n/2$ uses eight products of blocks. In 1969 Strassen observed that seven suffice~\cite{strassen1969}. With the blocks $A_{IJ}$ and $B_{JK}$, form
\[
\begin{aligned}
M_1&=(A_{11}+A_{22})(B_{11}+B_{22}),& M_2&=(A_{21}+A_{22})B_{11},\\
M_3&=A_{11}(B_{12}-B_{22}),& M_4&=A_{22}(B_{21}-B_{11}),\\
M_5&=(A_{11}+A_{12})B_{22},& M_6&=(A_{21}-A_{11})(B_{11}+B_{12}),\\
M_7&=(A_{12}-A_{22})(B_{21}+B_{22});&&
\end{aligned}
\]
then the blocks of the product are $C_{11}=M_1+M_4-M_5+M_7$, $C_{12}=M_3+M_5$, $C_{21}=M_2+M_4$ and $C_{22}=M_1-M_2+M_3+M_6$, as multiplying out confirms; the order of every product is kept, so the formulas hold for blocks, not only for numbers. Applying the same scheme recursively to the seven products of order $n/2$, the cost $T(n)$ satisfies $T(n)=7T(n/2)+O(n^2)$, whose solution is $T(n)=O(n^{\log_27})$, with $\log_27=2.8073\ldots$, instead of $O(n^3)$. The infimum $\omega$ of the exponents $\alpha$ for which matrices can be multiplied with $O(n^\alpha)$ operations is the \emph{exponent of matrix multiplication}. Clearly $\omega\ge2$, since the $n^2$ entries of the product must all be computed, and a long series of improvements has brought the upper bound below $2.372$; Alman, Duan, Vassilevska Williams, Xu, Xu and Zhou proved $\omega<2.371339$~\cite{alman2025}. Whether $\omega=2$ is open. Only Strassen's algorithm and its variants are used in practice, for large matrices, because the later algorithms win only for astronomically large $n$; Strassen's algorithm also satisfies a weaker bound on rounding errors than the ordinary product~\cite[Chap.~23]{higham2002accuracy}.
\end{note}

\section{Trace, the Frobenius Inner Product, and Matrix Norms}
\label{sec:trace-norms}

\subsection{The Trace}

\begin{definition}[Trace]
\label{def:trace}
The \term{trace} of a square matrix $A\in\mathbb R^{n\times n}$ is the sum of its diagonal entries, $\operatorname{tr}A\coloneqq a_{11}+a_{22}+\cdots+a_{nn}$.
\end{definition}

The trace of the identity $I_n$ is $n$, that of the second difference matrix $T_n$ is $2n$, and that of a rotation $R_\varphi$ is $2\cos\varphi$. Unlike the determinant, the trace is linear, and its most important property concerns products.

\begin{proposition}[Properties of the trace]
\label{prop:trace}
\begin{enumerate}
\item $\operatorname{tr}(sA+tB)=s\operatorname{tr}A+t\operatorname{tr}B$ and $\operatorname{tr}A^{\mathsf T}=\operatorname{tr}A$ for $A,B\in\mathbb R^{n\times n}$.
\item $\operatorname{tr}(AB)=\operatorname{tr}(BA)$ for $A\in\mathbb R^{m\times n}$ and $B\in\mathbb R^{n\times m}$; in particular $\operatorname{tr}(\mathbf u\mathbf v^{\mathsf T})=\langle\mathbf u,\mathbf v\rangle$.
\item $\operatorname{tr}(ABC)=\operatorname{tr}(BCA)=\operatorname{tr}(CAB)$ whenever the products are defined and square.
\item $\operatorname{tr}(S^{-1}AS)=\operatorname{tr}A$ for every invertible $S\in\mathbb R^{n\times n}$.
\end{enumerate}
\end{proposition}

\begin{proof}
(i) The diagonal entries of $sA+tB$ are $sa_{ii}+tb_{ii}$, and transposition does not move the diagonal. (ii) Exchanging the order of summation,
\[
\operatorname{tr}(AB)=\sum_{i=1}^m\sum_{j=1}^na_{ij}b_{ji}=\sum_{j=1}^n\sum_{i=1}^mb_{ji}a_{ij}=\operatorname{tr}(BA).
\]
For $A=\mathbf u$ and $B=\mathbf v^{\mathsf T}$, $\operatorname{tr}(\mathbf u\mathbf v^{\mathsf T})=\operatorname{tr}(\mathbf v^{\mathsf T}\mathbf u)=\langle\mathbf v,\mathbf u\rangle$. (iii) Apply (ii) to the two factors $A$ and $BC$, and then to $AB$ and $C$. (iv) By (iii), $\operatorname{tr}(S^{-1}AS)=\operatorname{tr}(ASS^{-1})=\operatorname{tr}A$.
\end{proof}

In (ii) the products $AB$ and $BA$ may have different sizes, $m\times m$ and $n\times n$, and yet they have the same trace. The cyclic rule (iii) allows the factors to be rotated, but not rearranged arbitrarily: with the matrices $E_{ij}\coloneqq\mathbf e_i\mathbf e_j^{\mathsf T}$, which have a single entry $1$ in position $(i,j)$, $\operatorname{tr}(E_{11}E_{12}E_{21})=\operatorname{tr}E_{11}=1$ but $\operatorname{tr}(E_{11}E_{21}E_{12})=\operatorname{tr}0=0$. Part (iv) says that the trace does not depend on the coordinates in which a linear map is described, a fact that \cref{ch:eigenvalues} explains: the trace is the sum of the eigenvalues. The trace of the orthogonal projection $\mathbf u\mathbf u^{\mathsf T}$ onto a line, with $\|\mathbf u\|=1$, is $\langle\mathbf u,\mathbf u\rangle=1$, the dimension of the line, and \cref{prop:projection} shows that the trace of every idempotent matrix, orthogonal projection or not, is the dimension of its image.

\subsection{The Frobenius Inner Product}

The set $\mathbb R^{m\times n}$ is a copy of $\mathbb R^{mn}$, and the Euclidean inner product of $\mathbb R^{mn}$ becomes, on matrices, the \term{Frobenius inner product}
\[
\langle A,B\rangle_{\mathrm F}\coloneqq\sum_{i=1}^m\sum_{j=1}^na_{ij}b_{ij}=\operatorname{tr}(A^{\mathsf T}B),
\]
where the last equality holds because the $j$th diagonal entry of $A^{\mathsf T}B$ is $\sum_ia_{ij}b_{ij}$. The corresponding norm is the Frobenius norm $\|A\|_{\mathrm F}=\langle A,A\rangle_{\mathrm F}^{1/2}=\bigl(\sum_{i,j}a_{ij}^2\bigr)^{1/2}$ of \cref{prop:matrix-norms}, which showed that $\|A\mathbf x\|\le\|A\|_{\mathrm F}\|\mathbf x\|$ and $\|AB\|_{\mathrm F}\le\|A\|_{\mathrm F}\|B\|_{\mathrm F}$. Everything proved about inner products on $\mathbb R^{mn}$ holds for it; the Cauchy--Schwarz inequality, for instance, reads $|\operatorname{tr}(A^{\mathsf T}B)|\le\|A\|_{\mathrm F}\|B\|_{\mathrm F}$. Two further facts make the Frobenius norm convenient: it adds up over columns, $\|A\|_{\mathrm F}^2=\sum_j\|\mathbf a_j\|^2$, and it is unchanged by orthogonal transformations, $\|QA\|_{\mathrm F}=\|A\|_{\mathrm F}=\|AQ'\|_{\mathrm F}$ for orthogonal $Q$ and $Q'$, because $Q$ preserves the length of every column of $A$, and $\|AQ'\|_{\mathrm F}=\|Q'^{\mathsf T}A^{\mathsf T}\|_{\mathrm F}$. Chapter~21 measures the distance of a matrix from the matrices of low rank in this norm.

\subsection{Operator Norms}

The operator norm $\|A\|_2$ of \cref{prop:matrix-norms} is the largest factor by which $A$ stretches a vector, measured in the Euclidean norm. Measuring in the other norms of \cref{def:norm} gives other matrix norms.

\begin{definition}[Operator norms]
\label{def:operator-norm}
For $A\in\mathbb R^{m\times n}$ and $p\in\{1,2,\infty\}$, the \term{operator norm} induced by $\|\cdot\|_p$ is
\[
\|A\|_p\coloneqq\sup\bigl\{\|A\mathbf x\|_p: \mathbf x\in\mathbb R^n,\ \|\mathbf x\|_p=1\bigr\}.
\]
\end{definition}

For $p=2$ this is the operator norm of \cref{prop:matrix-norms}. Unlike $\|A\|_2$, which Chapter~21 computes as the largest singular value of $A$, the norms for $p=1$ and $p=\infty$ have simple formulas.

\begin{theorem}[The \texorpdfstring{$1$}{1}-norm and the \texorpdfstring{$\infty$}{infinity}-norm of a matrix]
\label{thm:norm-formulas}
For $A\in\mathbb R^{m\times n}$,
\[
\|A\|_1=\max_{1\le j\le n}\sum_{i=1}^m|a_{ij}|,\qquad\|A\|_\infty=\max_{1\le i\le m}\sum_{j=1}^n|a_{ij}| ,
\]
the largest column sum and the largest row sum of the absolute values of the entries. In particular $\|A^{\mathsf T}\|_1=\|A\|_\infty$, and both suprema are attained.
\end{theorem}

\begin{proof}
Let $c$ be the largest column sum, attained in column $j^*$. For every $\mathbf x$,
\[
\|A\mathbf x\|_1=\sum_i\Bigl|\sum_ja_{ij}x_j\Bigr|\le\sum_i\sum_j|a_{ij}||x_j|=\sum_j|x_j|\sum_i|a_{ij}|\le c\,\|\mathbf x\|_1 ,
\]
so $\|A\|_1\le c$; and $\mathbf x=\mathbf e_{j^*}$, of $1$-norm $1$, gives $\|A\mathbf e_{j^*}\|_1=c$. Let $r$ be the largest row sum, attained in row $i^*$. For every $\mathbf x$ and every $i$,
\[
|(A\mathbf x)_i|\le\sum_j|a_{ij}||x_j|\le\|\mathbf x\|_\infty\sum_j|a_{ij}|\le r\,\|\mathbf x\|_\infty ,
\]
so $\|A\|_\infty\le r$. If $A=0$ there is nothing more to prove. Otherwise $r>0$, and the vector $\mathbf x$ with $x_j\coloneqq\operatorname{sign}a_{i^*j}$ has $\|\mathbf x\|_\infty=1$, since some $a_{i^*j}\ne0$, and $(A\mathbf x)_{i^*}=\sum_j|a_{i^*j}|=r$. The rows of $A^{\mathsf T}$ are the columns of $A$, which gives the last claim.
\end{proof}

\begin{proposition}[Operator norms]
\label{prop:operator-norms}
Let $p\in\{1,2,\infty\}$. For all matrices whose shapes fit:
\begin{enumerate}
\item $\|A\mathbf x\|_p\le\|A\|_p\|\mathbf x\|_p$, and $\|A\|_p$ is the smallest number with this property;
\item $\|\cdot\|_p$ is a norm on $\mathbb R^{m\times n}$, and $\|I\|_p=1$;
\item $\|AB\|_p\le\|A\|_p\|B\|_p$;
\item $|a_{ij}|\le\|A\|_p\le n\max_{k,l}|a_{kl}|$ for $A\in\mathbb R^{n\times n}$ and all $i,j$, and the same holds for $\|A\|_{\mathrm F}$.
\end{enumerate}
\end{proposition}

\begin{proof}
For $p=2$, (i) to (iii) are \cref{prop:matrix-norms}, and the proof given there works for $p=1$ and $p=\infty$ as well, since it uses only the norm properties of $\|\cdot\|_p$ and the finiteness of the supremum, which \cref{thm:norm-formulas} provides. $\|I\|_p=1$ because $\|I\mathbf x\|_p=\|\mathbf x\|_p$. For (iv), $|a_{ij}|$ is at most the $j$th column sum and the $i$th row sum, and at most $\|A\|_2$ by \cref{prop:matrix-norms}\,(iii); and a column or row sum of $n$ terms is at most $n\max|a_{kl}|$, as is $\|A\|_{\mathrm F}=(\sum a_{kl}^2)^{1/2}\le(n^2\max a_{kl}^2)^{1/2}$, which bounds $\|A\|_2$ by \cref{prop:matrix-norms}\,(i).
\end{proof}

For the matrix with the rows $(1,-2)$ and $(3,4)$, the column sums are $4$ and $6$ and the row sums $3$ and $7$, so $\|A\|_1=6$ and $\|A\|_\infty=7$, while $\|A\|_{\mathrm F}=\sqrt{30}=5.477\ldots$ and, as Chapter~21 will show, $\|A\|_2=(15+5\sqrt5)^{1/2}=5.1166\ldots$. The four numbers differ, but by (iv) they are within a factor $n$ of each other and of the largest entry. In particular, a sequence of matrices $A_k$ converges entrywise to $A$ if and only if $\|A_k-A\|\to0$ in any of these norms, as for vectors in \cref{prop:componentwise}. The submultiplicativity (iii) is what makes operator norms useful for products and powers: $\|A^k\|_p\le\|A\|_p^k$.

\subsection{The Neumann Series}

For a number $a$ with $|a|<1$, the geometric series gives $(1-a)^{-1}=\sum_ka^k$ (\cref{prop:geometric-series}). The same is true for matrices, with the absolute value replaced by an operator norm, and \cref{prop:nilpotent} was the special case of a nilpotent matrix, for which the series is finite.

\begin{theorem}[Neumann series]
\label{thm:neumann}
Let $A\in\mathbb R^{n\times n}$ and $\|A\|<1$, where $\|\cdot\|$ is one of the operator norms $\|\cdot\|_1$, $\|\cdot\|_2$, $\|\cdot\|_\infty$. Then $I-A$ is invertible, the series $\sum_{k=0}^\infty A^k$ converges entrywise, and
\[
(I-A)^{-1}=\sum_{k=0}^\infty A^k,\qquad\bigl\|(I-A)^{-1}\bigr\|\le\frac1{1-\|A\|},\qquad\Bigl\|(I-A)^{-1}-\sum_{k=0}^NA^k\Bigr\|\le\frac{\|A\|^{N+1}}{1-\|A\|}.
\]
\end{theorem}

\begin{proof}
Write $q\coloneqq\|A\|$ and $S_N\coloneqq\sum_{k=0}^NA^k$. By \cref{prop:operator-norms}, $|(A^k)_{ij}|\le\|A^k\|\le q^k$, so for every $i,j$ the series $\sum_k(A^k)_{ij}$ converges absolutely by comparison with the geometric series (\cref{thm:comparison-series,thm:absolute-series}). Hence $S_N$ converges entrywise to a matrix $S$. The sum telescopes,
\[
(I-A)S_N=S_N-AS_N=\sum_{k=0}^NA^k-\sum_{k=1}^{N+1}A^k=I-A^{N+1},
\]
and similarly $S_N(I-A)=I-A^{N+1}$. As $N\to\infty$, each entry of $(I-A)S_N$ is a fixed linear combination of entries of $S_N$ and converges to the corresponding entry of $(I-A)S$, while $|(A^{N+1})_{ij}|\le q^{N+1}\to0$. So $(I-A)S=I=S(I-A)$, and $S=(I-A)^{-1}$. Next, $(I-A)(S-S_N)=I-(I-A^{N+1})=A^{N+1}$, so $S-S_N=SA^{N+1}$. For the bounds, the triangle inequality and \cref{prop:operator-norms}\,(iii) and~(iv) give
\[
\|S\|\le\|S_N\|+\|S-S_N\|\le\sum_{k=0}^Nq^k+n\max_{i,j}|(S-S_N)_{ij}| ;
\]
letting $N\to\infty$ gives $\|S\|\le1/(1-q)$. Finally $\|S-S_N\|=\|SA^{N+1}\|\le\|S\|q^{N+1}\le q^{N+1}/(1-q)$.
\end{proof}

The condition $\|A\|<1$ is sufficient but not necessary for convergence. What decides is whether the powers $A^k$ tend to zero, and this happens exactly when all eigenvalues of $A$ have modulus less than $1$, a fact that \cref{ch:eigenvalues} proves. \Cref{fig:neumann} shows the difference. For the symmetric matrix $A_1$ with the rows $(1/2,1/4)$ and $(1/4,1/2)$, $\|A_1\|_\infty=3/4$, and the norms of the powers decrease like $(3/4)^k$. The matrix $A_2$ with the rows $(4/5,2)$ and $(0,4/5)$ has $\|A_2\|_\infty=14/5>1$, and the theorem does not apply; yet $A_2^k$ has the rows $((4/5)^k,\,2k(4/5)^{k-1})$ and $(0,(4/5)^k)$, which tend to zero, since the geometric factor eventually beats the linear growth of $k$. The norms of the powers first grow, to $4.5$ at $k=4$, and then decay, and the Neumann series still converges, to the inverse of $I-A_2$, which has the rows $(5,50)$ and $(0,5)$. Such \emph{transient growth} is typical of matrices that are far from symmetric, or more precisely far from normal in the sense of \cref{sec:complex-matrices}, and it has practical consequences for iterative methods (Volume~IX).

\begin{figure}
\sdfloat{%
\begin{tikzpicture}
\begin{axis}[sdplot, width=0.6\textwidth, height=0.4\textwidth, ymode=log,
  xmin=0, xmax=60, ymin=1e-8, ymax=30, xtick={0,10,20,30,40,50,60},
  ytick={1e-8,1e-6,1e-4,1e-2,1}, xlabel={$k$}, legend pos=north east]
\addplot[black, line width=0.7pt, domain=0:60, samples=61] {0.75^x};
\addlegendentry{$\|A_1^k\|_\infty$}
\addplot[black, dashed, line width=0.7pt, domain=0:60, samples=61] {0.8^x*(1+2.5*x)};
\addlegendentry{$\|A_2^k\|_\infty$}
\end{axis}
\end{tikzpicture}}{The norms of the powers of two matrices whose Neumann series converge. For $A_1$, with $\|A_1\|_\infty=3/4$, they decrease geometrically, as \cref{thm:neumann} guarantees. For $A_2$, with $\|A_2\|_\infty=14/5$, they first grow to $4.5$ and then decay like $k(4/5)^k$; the theorem does not apply, but the series converges all the same. The values at integer $k$ are joined by lines.\label{fig:neumann}}
\end{figure}

The theorem has an important consequence: an invertible matrix remains invertible under small perturbations, and its inverse changes little.

\begin{corollary}[Perturbation of the inverse]
\label{cor:perturbed-inverse}
Let $A\in\mathbb R^{n\times n}$ be invertible, $E\in\mathbb R^{n\times n}$, and $\|A^{-1}\|\|E\|<1$ in one of the norms of \cref{thm:neumann}. Then $A+E$ is invertible, and
\[
\bigl\|(A+E)^{-1}\bigr\|\le\frac{\|A^{-1}\|}{1-\|A^{-1}\|\|E\|},\qquad\bigl\|(A+E)^{-1}-A^{-1}\bigr\|\le\frac{\|A^{-1}\|^2\|E\|}{1-\|A^{-1}\|\|E\|}.
\]
\end{corollary}

\begin{proof}
Write $A+E=A(I-F)$ with $F\coloneqq-A^{-1}E$, so that $\|F\|\le\|A^{-1}\|\|E\|<1$. By \cref{thm:neumann}, $I-F$ is invertible with $\|(I-F)^{-1}\|\le1/(1-\|F\|)$, so $A+E$ is invertible with the inverse $(I-F)^{-1}A^{-1}$, which gives the first bound. For the second, $(A+E)^{-1}-A^{-1}=-(A+E)^{-1}EA^{-1}$, as multiplying by $A+E$ from the left confirms, and the norm of the right side is at most $\|(A+E)^{-1}\|\|E\|\|A^{-1}\|$.
\end{proof}

By the corollary, every matrix close enough to an invertible matrix is invertible: the invertible matrices form an open set in $\mathbb R^{n\times n}$, and the inverse depends continuously on the matrix, as \cref{thm:invertible-matrices}\,(iii) also stated. Dividing the second bound by $\|A^{-1}\|$ gives a bound for the relative change of the inverse,
\[
\frac{\|(A+E)^{-1}-A^{-1}\|}{\|A^{-1}\|}\le\frac{\kappa(A)\,\|E\|/\|A\|}{1-\kappa(A)\,\|E\|/\|A\|},\qquad\kappa(A)\coloneqq\|A\|\,\|A^{-1}\|,
\]
in which the relative perturbation $\|E\|/\|A\|$ of the matrix is amplified by the \term{condition number} $\kappa(A)$. This number, the matrix counterpart of the condition number of \cref{sec:conditioning}, governs the accuracy of every method for linear systems; Chapter~21 and Volume~VIII study it.

\section{Complex Vectors and Matrices}
\label{sec:complex-matrices}

\subsection{Complex Entries and the Conjugate Transpose}

Vectors and matrices with complex entries arise as soon as eigenvalues are complex (\cref{ch:eigenvalues}), and they are the natural language of the Fourier transform (Volume~XI). The set $\mathbb C^n$ of $n$-tuples of complex numbers, with addition and multiplication by complex scalars component by component, satisfies the eight rules of \cref{prop:vector-rules}, and $\mathbb C^{m\times n}$ denotes the complex $m\times n$ matrices. Everything in \cref{sec:matrices-linear-maps,sec:matrix-multiplication,sec:inverse,sec:block-matrices} holds verbatim with complex numbers in place of real ones, since only the rules of arithmetic of a field were used; this includes \cref{thm:one-sided-square}, whose proof by elimination in \cref{ch:linear-systems} works over $\mathbb C$ as well. One thing must change: the inner product.

The \term{conjugate} $\bar A$ of a complex matrix has the entries $\bar a_{ij}$, and the \term{conjugate transpose} is $A^{\mathsf H}\coloneqq\bar A^{\mathsf T}$, with the entries $(A^{\mathsf H})_{ji}=\bar a_{ij}$. For a real matrix, $A^{\mathsf H}=A^{\mathsf T}$. For example,
\[
\begin{bmatrix}1+i&2\\3i&4-i\end{bmatrix}^{\mathsf H}=\begin{bmatrix}1-i&-3i\\2&4+i\end{bmatrix}.
\]
Many authors write $A^*$ instead of $A^{\mathsf H}$; we keep the star for other uses, such as the optimal value $t^*$ of \cref{prop:projection-line}, and the letter H recalls Hermite, after whom the matrices with $A^{\mathsf H}=A$ are named.

\begin{proposition}[Rules for the conjugate transpose]
\label{prop:conjugate-transpose}
Whenever the shapes fit, $(A^{\mathsf H})^{\mathsf H}=A$, $(A+B)^{\mathsf H}=A^{\mathsf H}+B^{\mathsf H}$, $(tA)^{\mathsf H}=\bar tA^{\mathsf H}$ for $t\in\mathbb C$, $(AB)^{\mathsf H}=B^{\mathsf H}A^{\mathsf H}$, and $(A^{-1})^{\mathsf H}=(A^{\mathsf H})^{-1}$ for invertible $A$.
\end{proposition}

\begin{proof}
Conjugation of complex numbers respects sums and products and is its own inverse (\cref{prop:conjugate}\,(i)), so $\overline{AB}=\bar A\bar B$ entrywise, and the rules follow from those for the transpose (\cref{prop:transpose-rules,prop:inverse-rules}): for instance $(AB)^{\mathsf H}=(\bar A\bar B)^{\mathsf T}=\bar B^{\mathsf T}\bar A^{\mathsf T}=B^{\mathsf H}A^{\mathsf H}$.
\end{proof}

\subsection{The Hermitian Inner Product}

The formula $\sum_ix_iy_i$ does not give a useful inner product on $\mathbb C^n$: for $\mathbf x=(1,i)$ it gives $1+i^2=0$, so a nonzero vector would have length zero. The remedy is the conjugate, as for the modulus of a complex number, $|z|^2=\bar zz$.

\begin{definition}[Hermitian inner product]
\label{def:hermitian-inner-product}
The \term{inner product} of $\mathbf x,\mathbf y\in\mathbb C^n$ is
\[
\langle\mathbf x,\mathbf y\rangle\coloneqq\mathbf x^{\mathsf H}\mathbf y=\sum_{i=1}^n\bar x_iy_i ,
\]
and the \term{norm} of $\mathbf x$ is $\|\mathbf x\|\coloneqq\langle\mathbf x,\mathbf x\rangle^{1/2}=\bigl(\sum_i|x_i|^2\bigr)^{1/2}$. For real vectors these are the inner product and the norm of \cref{def:inner-product}.
\end{definition}

The inner product is conjugated in its \emph{first} argument and linear in its second. This is the convention of numerical linear algebra and of physics, and it matches the matrix notation $\mathbf x^{\mathsf H}\mathbf y$; many texts in pure mathematics, Axler's among them~\cite[Chap.~6]{axler2024}, make the inner product linear in its first argument and conjugate the second instead. The two conventions give complex conjugate values, and formulas taken from other sources must be adjusted accordingly.

\begin{proposition}[Properties of the Hermitian inner product]
\label{prop:hermitian-inner-product}
For all $\mathbf x,\mathbf y,\mathbf z\in\mathbb C^n$, $s,t\in\mathbb C$ and $A\in\mathbb C^{m\times n}$:
\begin{enumerate}
\item $\langle\mathbf y,\mathbf x\rangle=\overline{\langle\mathbf x,\mathbf y\rangle}$;
\item $\langle\mathbf z,s\mathbf x+t\mathbf y\rangle=s\langle\mathbf z,\mathbf x\rangle+t\langle\mathbf z,\mathbf y\rangle$ and $\langle s\mathbf x+t\mathbf y,\mathbf z\rangle=\bar s\langle\mathbf x,\mathbf z\rangle+\bar t\langle\mathbf y,\mathbf z\rangle$;
\item $\langle\mathbf x,\mathbf x\rangle$ is real and positive for $\mathbf x\ne\mathbf 0$;
\item $|\langle\mathbf x,\mathbf y\rangle|\le\|\mathbf x\|\|\mathbf y\|$ (Cauchy--Schwarz), and $\|\cdot\|$ is a norm on $\mathbb C^n$;
\item $\langle A\mathbf x,\mathbf w\rangle=\langle\mathbf x,A^{\mathsf H}\mathbf w\rangle$ for $\mathbf w\in\mathbb C^m$.
\end{enumerate}
\end{proposition}

\begin{proof}
(i) $\overline{\sum_i\bar x_iy_i}=\sum_ix_i\bar y_i$. (ii) The first rule holds termwise, and the second follows from the first and (i). (iii) $\bar x_ix_i=|x_i|^2$. (iv) Let $\mathbf y\ne\mathbf 0$ (otherwise both sides vanish) and $t\coloneqq\langle\mathbf y,\mathbf x\rangle/\|\mathbf y\|^2$. By (i) and (ii),
\[
0\le\|\mathbf x-t\mathbf y\|^2=\|\mathbf x\|^2-t\langle\mathbf x,\mathbf y\rangle-\bar t\langle\mathbf y,\mathbf x\rangle+|t|^2\|\mathbf y\|^2=\|\mathbf x\|^2-\frac{|\langle\mathbf x,\mathbf y\rangle|^2}{\|\mathbf y\|^2},
\]
since each of the last three terms equals $|\langle\mathbf x,\mathbf y\rangle|^2/\|\mathbf y\|^2$ in absolute value, with the signs $-,-,+$. For the triangle inequality, $\|\mathbf x+\mathbf y\|^2=\|\mathbf x\|^2+2\operatorname{Re}\langle\mathbf x,\mathbf y\rangle+\|\mathbf y\|^2\le(\|\mathbf x\|+\|\mathbf y\|)^2$, because $\operatorname{Re}z\le|z|$; positivity and homogeneity are clear. (v) $\langle A\mathbf x,\mathbf w\rangle=(A\mathbf x)^{\mathsf H}\mathbf w=\mathbf x^{\mathsf H}A^{\mathsf H}\mathbf w$.
\end{proof}

The norms $\|\mathbf x\|_1=\sum_i|x_i|$ and $\|\mathbf x\|_\infty=\max_i|x_i|$ are defined on $\mathbb C^n$ with the modulus in place of the absolute value, and the results of \cref{sec:trace-norms} carry over, with $A^{\mathsf H}$ in place of $A^{\mathsf T}$ where the transpose meets an inner product: $\langle A,B\rangle_{\mathrm F}=\operatorname{tr}(A^{\mathsf H}B)$, and in the proof of \cref{thm:norm-formulas} the signs of the entries are replaced by the complex numbers $\bar a_{i^*j}/|a_{i^*j}|$ of modulus one.

\subsection{Hermitian, Unitary and Normal Matrices}

The real notions of \cref{sec:transpose,sec:special-matrices} that involve the transpose have complex counterparts with the conjugate transpose.

\begin{definition}[Hermitian, unitary and normal matrices]
\label{def:hermitian-unitary}
A matrix $A\in\mathbb C^{n\times n}$ is \term{Hermitian} if $A^{\mathsf H}=A$, \term{skew-Hermitian} if $A^{\mathsf H}=-A$, \term{unitary} if $A^{\mathsf H}A=I$, and \term{normal} if $A^{\mathsf H}A=AA^{\mathsf H}$.
\end{definition}

\begin{proposition}
\label{prop:hermitian-unitary}
Let $A,U\in\mathbb C^{n\times n}$.
\begin{enumerate}
\item If $A$ is Hermitian, its diagonal entries are real, and $\langle\mathbf x,A\mathbf x\rangle$ is real for every $\mathbf x\in\mathbb C^n$.
\item The following are equivalent: $U$ is unitary; the columns of $U$ are orthonormal; $\langle U\mathbf x,U\mathbf y\rangle=\langle\mathbf x,\mathbf y\rangle$ for all $\mathbf x,\mathbf y$; $\|U\mathbf x\|=\|\mathbf x\|$ for all $\mathbf x$. A unitary matrix is invertible with $U^{-1}=U^{\mathsf H}$.
\item Hermitian, skew-Hermitian and unitary matrices are normal.
\end{enumerate}
\end{proposition}

\begin{proof}
(i) $a_{ii}=\bar a_{ii}$, and $\overline{\langle\mathbf x,A\mathbf x\rangle}=\langle A\mathbf x,\mathbf x\rangle=\langle\mathbf x,A^{\mathsf H}\mathbf x\rangle=\langle\mathbf x,A\mathbf x\rangle$ by \cref{prop:hermitian-inner-product}\,(i) and (v). (ii) The proof of \cref{thm:orthogonal} goes through with one change in the step from norms to inner products: by \cref{prop:hermitian-inner-product}, $\operatorname{Re}\langle\mathbf x,\mathbf y\rangle=(\|\mathbf x+\mathbf y\|^2-\|\mathbf x-\mathbf y\|^2)/4$, and $\operatorname{Im}\langle\mathbf x,\mathbf y\rangle=\operatorname{Re}\langle\mathbf x,-i\mathbf y\rangle$, so norms determine inner products again. (iii) For Hermitian $A$, both $A^{\mathsf H}A$ and $AA^{\mathsf H}$ equal $A^2$, and for skew-Hermitian $A$ both equal $-A^2$; for unitary $U$, $U^{\mathsf H}U=I=UU^{\mathsf H}$ by (ii).
\end{proof}

The matrix with the rows $(2,i)$ and $(-i,3)$ is Hermitian, and for $\mathbf x=(1,1)$ the number $\langle\mathbf x,A\mathbf x\rangle=(2+i)+(3-i)=5$ is indeed real. The matrix
\[
U=\frac1{\sqrt2}\begin{bmatrix}1&1\\i&-i\end{bmatrix}
\]
is unitary: its columns have length $1$, and their inner product is $(1\cdot1+\bar i\cdot(-i))/2=(1+i^2)/2=0$. The most important unitary matrix is that of the discrete Fourier transform. For $n\ge1$ and $\omega=e^{2\pi i/n}$, let $F_n$ be the $n\times n$ matrix with the entries $\omega^{jk}$, $j,k=0,\ldots,n-1$. The $(j,k)$ entry of $F_n^{\mathsf H}F_n$ is $\sum_l\bar\omega^{lj}\omega^{lk}=\sum_l\omega^{l(k-j)}$, since $\bar\omega=\omega^{-1}$, and by \cref{prop:roots-of-unity} this is $n$ for $j=k$ and $0$ otherwise. So $F_n/\sqrt n$ is unitary, and the inverse of $F_n$ is $F_n^{\mathsf H}/n$; Volume~XI builds the fast Fourier transform on this fact, with the sign convention of the exponent discussed there.

Real orthogonal matrices are unitary, real symmetric matrices are Hermitian, and all of them are normal. A normal matrix need not be of any of these kinds: the matrix with the rows $(1,-1)$ and $(1,1)$, which is $\sqrt2R_{\pi/4}$, is normal but neither Hermitian nor unitary. The shear with the rows $(1,1)$ and $(0,1)$ is not normal, since $A^{\mathsf H}A$ has the rows $(1,1)$, $(1,2)$ and $AA^{\mathsf H}$ the rows $(2,1)$, $(1,1)$. Normality looks like a technical condition, but it characterises the matrices that can be diagonalised by a unitary matrix, as Chapter~22 shows, and the eigenvalues of Hermitian matrices are real, as \cref{ch:symmetric-matrices} shows. \Cref{tab:real-complex} collects the correspondence between the real and the complex notions.

\begin{table}
\sdfloat{%
\begin{sdtable}{colspec={ll}, row{1}={font=\small\bfseries}}
\toprule
real & complex \\
\midrule
$\mathbb R^n$, $\mathbb R^{m\times n}$ & $\mathbb C^n$, $\mathbb C^{m\times n}$ \\
transpose $A^{\mathsf T}$ & conjugate transpose $A^{\mathsf H}=\bar A^{\mathsf T}$ \\
$\langle\mathbf x,\mathbf y\rangle=\mathbf x^{\mathsf T}\mathbf y=\sum x_iy_i$ & $\langle\mathbf x,\mathbf y\rangle=\mathbf x^{\mathsf H}\mathbf y=\sum\bar x_iy_i$ \\
$\|\mathbf x\|^2=\sum x_i^2$ & $\|\mathbf x\|^2=\sum|x_i|^2$ \\
symmetric, $A^{\mathsf T}=A$ & Hermitian, $A^{\mathsf H}=A$ \\
antisymmetric, $A^{\mathsf T}=-A$ & skew-Hermitian, $A^{\mathsf H}=-A$ \\
orthogonal, $Q^{\mathsf T}Q=I$ & unitary, $U^{\mathsf H}U=I$ \\
normal, $A^{\mathsf T}A=AA^{\mathsf T}$ & normal, $A^{\mathsf H}A=AA^{\mathsf H}$ \\
quadratic form $\langle\mathbf x,A\mathbf x\rangle$, $A$ symmetric & $\langle\mathbf x,A\mathbf x\rangle$ real for Hermitian $A$ \\
$\langle A,B\rangle_{\mathrm F}=\operatorname{tr}(A^{\mathsf T}B)$ & $\langle A,B\rangle_{\mathrm F}=\operatorname{tr}(A^{\mathsf H}B)$ \\
\bottomrule
\end{sdtable}}{Real notions and their complex counterparts. Every real notion is the special case of the complex one for real entries; in each formula the transpose becomes the conjugate transpose.\label{tab:real-complex}}
\end{table}

Complex numbers themselves can be represented by real matrices. Multiplication by $z=a+bi$ maps $x+iy$ to $(ax-by)+i(bx+ay)$, a real linear map of the plane with the matrix that has the rows $(a,-b)$ and $(b,a)$; writing $z=re^{i\theta}$ (\cref{sec:complex-exp}), this matrix is $rR_\theta$, a rotation followed by a scaling. Sums and products of complex numbers correspond to sums and products of these matrices, and the conjugate $\bar z$ to the transpose. In this way a complex $n\times n$ matrix can be stored as a real $2n\times2n$ matrix, in which every entry appears twice. This doubles the storage, and a product of two such real matrices costs twice the arithmetic of the complex product, $(2n)^3=8n^3$ real multiplications against $n^3$ complex ones of four real multiplications each; this is why libraries work with complex numbers directly.

\section{Computing with Vectors and Matrices: A First Look}
\label{sec:computing-matrices}

\subsection{Counting Operations}

The cost of a numerical computation is traditionally measured by the number of arithmetic operations on floating-point numbers, \term{flops}, that it performs. The count is a crude measure, since the time a computer takes depends as much on the movement of data as on arithmetic, but it is easy to determine, independent of the machine, and it describes how the cost grows with the size of the problem. \Cref{tab:op-counts} lists the counts for the basic operations of this chapter.

\begin{table}
\sdfloat{%
\begin{sdtable}{colspec={lll}, row{1}={font=\small\bfseries}}
\toprule
operation & multiplications & additions \\
\midrule
$\langle\mathbf x,\mathbf y\rangle$, $\mathbf x,\mathbf y\in\mathbb R^n$ & $n$ & $n-1$ \\
$\mathbf y+t\mathbf x$, $\mathbf x,\mathbf y\in\mathbb R^n$ & $n$ & $n$ \\
$\mathbf u\mathbf v^{\mathsf T}$, $\mathbf u\in\mathbb R^m$, $\mathbf v\in\mathbb R^n$ & $mn$ & $0$ \\
$A\mathbf x$, $A\in\mathbb R^{m\times n}$ & $mn$ & $m(n-1)$ \\
$AB$, $A\in\mathbb R^{m\times n}$, $B\in\mathbb R^{n\times p}$ & $mnp$ & $mp(n-1)$ \\
$D\mathbf x$, $D$ diagonal of order $n$ & $n$ & $0$ \\
$U\mathbf x$, $U$ upper triangular of order $n$ & $n(n+1)/2$ & $n(n-1)/2$ \\
$UV$, $U,V$ upper triangular of order $n$ & $n(n+1)(n+2)/6$ & $(n-1)n(n+1)/6$ \\
$P_\sigma\mathbf x$, $P_\sigma$ a permutation matrix & $0$ & $0$ \\
$A\mathbf x$, $A$ of order $n$ with bandwidths $(p,q)$ & at most $(p+q+1)n$ & at most $(p+q)n$ \\
$H\mathbf x=\mathbf x-2\mathbf u\langle\mathbf u,\mathbf x\rangle$ & $2n+1$ & $2n-1$ \\
\bottomrule
\end{sdtable}}{Operation counts for basic operations with vectors and matrices, by the obvious algorithms. The product of two upper triangular matrices needs about one sixth of the multiplications of a general product, since the entry $(i,k)$ is a sum over $i\le j\le k$ only.\label{tab:op-counts}}
\end{table}

The counts fall into three levels. Operations on vectors, such as inner products and the update $\mathbf y+t\mathbf x$, cost $O(n)$ operations on $O(n)$ data; matrix--vector products cost $O(n^2)$ operations on $O(n^2)$ data; and matrix--matrix products cost $O(n^3)$ operations on $O(n^2)$ data. The standard interface for these operations, the Basic Linear Algebra Subprograms (BLAS), is organised in the same three levels~\cite{dongarra1990}, and the third level is the one that runs fastest on modern computers, because each number read from memory takes part in about $n$ operations, which hides the cost of moving data. Structure reduces the counts, often dramatically: a diagonal matrix acts in $n$ operations, a triangular one in about $n^2$ instead of $2n^2$, a banded one in proportion to its nonzero entries, a permutation in none at all, and a Householder reflection, kept as the vector $\mathbf u$, in about $4n$ instead of $2n^2$. The count for the product of two upper triangular matrices follows by grouping the entries $(i,k)$ with $i\le k$ by $d\coloneqq k-i$: there are $n-d$ of them, each a sum of $d+1$ products, and
\[
\sum_{d=0}^{n-1}(n-d)(d+1)=\frac{n(n+1)(n+2)}6 .
\]

\subsection{The Order of Evaluation}

Associativity says that $(AB)C=A(BC)$, but the two sides can cost very different amounts. For square matrices of order $n$ and a vector $\mathbf x$, computing $(AB)\mathbf x$ costs $n^3+n^2$ multiplications and $A(B\mathbf x)$ costs $2n^2$; for $n=1000$ the first is about $500$ times as expensive. Likewise $(\mathbf u\mathbf v^{\mathsf T})\mathbf x$ costs $2n^2$ multiplications and $\mathbf u(\mathbf v^{\mathsf T}\mathbf x)$ costs $2n$. In a chain of matrices of different shapes the effect is less obvious. For $A\in\mathbb R^{10\times100}$, $B\in\mathbb R^{100\times5}$ and $C\in\mathbb R^{5\times50}$, the grouping $(AB)C$ costs $10\cdot100\cdot5+10\cdot5\cdot50=7500$ multiplications and $A(BC)$ costs $100\cdot5\cdot50+10\cdot100\cdot50=75\,000$, ten times as many. The best grouping of a long chain can be found by dynamic programming~\cite[Sect.~14.2]{cormen2022}, a technique of Volume~IV. The rule of thumb is to multiply by vectors as early as possible and never to form a matrix that is only going to be applied to a vector.

The same rule applies to polynomials of matrices. To compute $p(A)\mathbf x$ for $p(x)=c_0+c_1x+\cdots+c_dx^d$, forming the powers of $A$ costs $d-1$ matrix products, about $dn^3$ multiplications. Horner's scheme (\cref{sec:positional}) in the nested form
\[
p(A)\mathbf x=c_0\mathbf x+A\Bigl(c_1\mathbf x+A\bigl(c_2\mathbf x+\cdots+A(c_d\mathbf x)\cdots\bigr)\Bigr)
\]
needs only $d$ matrix--vector products, about $dn^2$ multiplications, and the vectors $\mathbf y_d\coloneqq c_d\mathbf x$, $\mathbf y_k\coloneqq c_k\mathbf x+A\mathbf y_{k+1}$ for $k=d-1,\ldots,0$ end with $\mathbf y_0=p(A)\mathbf x$. Many algorithms of later volumes, among them the Krylov methods for large linear systems and eigenvalue problems (Volume~IX), use $A$ only through products $A\mathbf x$ and never form anything else.

\subsection{Solving Systems without Inverses}

The solution of $A\mathbf x=\mathbf b$ is $A^{-1}\mathbf b$, but it should almost never be computed by forming $A^{-1}$ and multiplying. \Cref{ch:linear-systems} solves the system by Gaussian elimination, which factors $A=LU$ into triangular matrices with about $2n^3/3$ flops and then solves two triangular systems with about $2n^2$ flops. Forming $A^{-1}$ costs about $2n^3$ flops, three times as much, and the subsequent product another $2n^2$. The inverse is also less accurate in floating-point arithmetic, and it destroys sparsity: the inverse of the tridiagonal matrix $T_n$ is full (\cref{prop:second-difference-inverse}), while $T_n$ itself can be factored and solved in a number of operations proportional to $n$. Higham discusses when an explicit inverse is nevertheless justified~\cite[Chap.~14]{higham2002accuracy}. The rule is therefore: whenever a formula contains $A^{-1}\mathbf b$ or $A^{-1}B$, read it as ``the solution of a linear system''. This applies to the formulas of this chapter as well, to the Schur complement $D-CA^{-1}B$, in which $A^{-1}B$ is computed by solving systems with the columns of $B$ as right-hand sides, and to the Sherman--Morrison formula, which needs only $A^{-1}\mathbf b$ and $A^{-1}\mathbf u$.

In floating-point arithmetic, every operation of this chapter is rounded, and even an inner product $\sum_ix_iy_i$ is computed only approximately, with an error that depends on the order of the additions; Volume~V analyses these errors, and Volume~VIII the errors of matrix computations.

\subsection{How the Time Grows}

The operation count predicts that the time of an $n\times n$ matrix product grows like $n^3$: doubling $n$ should multiply the time by eight. \Cref{tab:matmul-timing} shows a small experiment that tests the prediction, with two very different implementations. The first is \cref{alg:matmul} written as three nested loops in the interpreted language Python, with the matrices stored as lists of rows. The second is the matrix product of the stardust library, the subject of Volume~VI, which hands the work to a tuned implementation of the BLAS routine for matrix products (on this machine Apple's Accelerate library). The measurements were made on an Apple M4 Pro processor, with Python~3.15 and the version of stardust of 25 November 2026, and each time is the best of several repetitions, for random matrices with entries in $[0,1)$.

\begin{table}
\sdfloat{%
\begin{sdtable}{colspec={rllll}, row{1}={font=\small\bfseries}}
\toprule
$n$ & loops (s) & ratio & library (s) & ratio \\
\midrule
32 & $6.6\times10^{-4}$ & & $7.9\times10^{-7}$ & \\
64 & $5.3\times10^{-3}$ & 8.0 & $2.1\times10^{-6}$ & 2.6 \\
128 & $5.0\times10^{-2}$ & 9.5 & $1.0\times10^{-5}$ & 5.0 \\
256 & $4.3\times10^{-1}$ & 8.5 & $8.2\times10^{-5}$ & 7.9 \\
512 & $3.8$ & 8.9 & $3.2\times10^{-4}$ & 4.0 \\
1024 & & & $2.6\times10^{-3}$ & 7.9 \\
2048 & & & $2.1\times10^{-2}$ & 8.1 \\
4096 & & & $1.7\times10^{-1}$ & 8.0 \\
\bottomrule
\end{sdtable}}{The time in seconds of an $n\times n$ matrix product by three nested loops in Python and by a tuned library routine, and the ratio of each time to the one in the row above. The operation count predicts a ratio of $8$. The loops were not run beyond $n=512$, where one product takes almost four seconds.\label{tab:matmul-timing}}
\end{table}

Both columns confirm the cubic law where it should hold, and the deviations are instructive. The loops follow it from the start, with ratios between $8$ and $9.5$. The library routine is more than $10\,000$ times faster than the loops at $n=512$, although both perform the same $n^3$ multiplications: it uses the vector instructions and the cache hierarchy of the processor, as described in \cref{sec:matrix-multiplication}, while every operation of the loops passes through the interpreter. For small $n$ the library's ratios are well below $8$. At $n=32$ the product takes about ten times as long as $32^3$ multiplications would take at the rate reached for large $n$, so the time is dominated by work that does not grow like $n^3$, such as the fixed cost of a call. Between $n=256$ and $n=512$ the ratio drops to about $4$, and this is where the routine starts to run on several processor cores at once: in a repetition of the experiment, the processor time it consumed was equal to the elapsed time up to $n=256$, about $1.6$ times the elapsed time at $n=512$ and about $3$ times from $n=1024$ on. From $n=1024$ on, the ratio is $8$ to within the accuracy of the measurement, and the time per multiplication settles at about $2.4\times10^{-12}$ seconds. The cubic law describes the growth, not the constant, and the constant can differ by four orders of magnitude.

\bigskip
A matrix is a linear map written down in coordinates, and the matrix product is the composition of maps. This chapter has derived the arithmetic of matrices from that idea: the rules that hold, the commutative law and the cancellation law that fail, the transpose as the adjoint with respect to the inner product, and the inverse as the undoing of a map, which exists exactly when no information is lost. Special kinds of matrices, diagonal, triangular, permutation, orthogonal, banded, carry structure that makes them cheap and reliable, and block matrices organise large computations and lead to the Schur complement and the Sherman--Morrison--Woodbury formula. Norms measure matrices, and the Neumann series shows that invertibility survives small perturbations, with the condition number as the measure of sensitivity. The questions that this chapter has deliberately left open are the subject of the following chapters: when a system $A\mathbf x=\mathbf b$ has a solution and how to compute it (\cref{ch:linear-systems}), what the dimension of a span is (\cref{ch:vector-spaces,ch:linear-maps}), how determinants measure volume and decide invertibility (\cref{ch:determinants}), how to project onto a subspace and orthogonalise a family of vectors (\cref{ch:inner-products}), and which vectors a matrix merely stretches (\cref{ch:eigenvalues}).

\chapternotes

The word \emph{matrix} was introduced by Sylvester in 1850 for a rectangular array from which determinants can be formed~\cite{sylvester1850}, and Cayley's memoir of 1858~\cite{cayley1858} treated matrices for the first time as objects of an algebra in their own right, with addition, multiplication, the inverse and the observation that multiplication is not commutative. Linear substitutions and their composition had been studied long before, by Gauss and Eisenstein among others~\cite{katz2009}, and the vector notation of Gibbs and Heaviside, with its inner and cross products, has its own history~\cite{crowe1967}. Katz~\cite{katz2009} gives an account of the development of linear algebra as a whole.

Among introductory textbooks, Strang~\cite{strang2016} emphasises the column picture of $A\mathbf x$ and the four readings of the product, Axler~\cite{axler2024} develops linear algebra from linear maps rather than matrices, and Meyer~\cite{meyer2000} combines the theory with numerical and applied aspects. Horn and Johnson~\cite{horn2013} is the standard reference for matrix analysis, including partitioned matrices, norms and the Neumann series (their Chapters~0 and~5). For the computational side, Trefethen and Bau~\cite{trefethen1997} is a concise introduction, and Golub and Van Loan~\cite{golub2013} and Higham~\cite{higham2002accuracy} are the standard references for algorithms and their rounding errors; Golub and Van Loan discuss the loop orderings of the matrix product, which Dongarra, Gustavson and Karp analysed for vector computers~\cite{dongarra1984}, and Goto and van de Geijn describe how modern libraries reach the speed of \cref{tab:matmul-timing}~\cite{goto2008}. The Householder reflection is from~\cite{householder1958}; the Sherman--Morrison and Woodbury formulas from~\cite{sherman1950,woodbury1950}, with history in~\cite{henderson1981} and applications in~\cite{hager1989}; the Schur complement from~\cite{schur1917}, surveyed in~\cite{zhang2005schur}; Strassen's algorithm from~\cite{strassen1969}, and the bound on the exponent of matrix multiplication quoted in the note from~\cite{alman2025}. Lagrange's identity and many relatives of the Cauchy--Schwarz inequality are in Steele~\cite{steele2004}, and the inequalities behind the $p$-norms in Hardy, Littlewood and P\'olya~\cite{hardy1952inequalities}.

\chapterreferences
